\documentclass[11pt,a4paper]{article}

\usepackage[utf8]{inputenc}
\usepackage[T1]{fontenc}
\usepackage[english]{babel}
\usepackage{amsmath, amssymb, amsthm, mathtools, mathrsfs, bm}
\usepackage{geometry}
\usepackage{hyperref}
\usepackage{booktabs}
\usepackage{array}
\usepackage{lmodern}
\usepackage{graphicx}
\graphicspath{{./figs/}{./}}
\usepackage{float}
\usepackage{subcaption}
\usepackage{xcolor}
\usepackage{listings}
\usepackage{caption}
\captionsetup{font=small,labelfont=bf}

\geometry{margin=2.5cm}

\hypersetup{colorlinks=true, linkcolor=blue!50!black,
  citecolor=blue!50!black, urlcolor=blue!50!black}

\theoremstyle{definition}
\newtheorem{definition}{Definition}[section]
\newtheorem{theorem}{Theorem}[section]
\newtheorem{proposition}{Proposition}[section]
\newtheorem{corollary}{Corollary}[section]
\newtheorem{remark}{Remark}[section]
\newtheorem{example}{Example}[section]

\definecolor{codegray}{rgb}{0.96,0.96,0.96}
\definecolor{leankw}{rgb}{0.5,0.0,0.3}
\definecolor{leancomment}{rgb}{0.4,0.5,0.4}
\lstdefinelanguage{lean}{
  morekeywords={theorem,lemma,def,let,fun,by,exact,calc,rw,simp,ring,intros,intro,have,import,namespace,end,open,sorry,noncomputable,structure,axiom,forall,exists},
  sensitive=true,
  morecomment=[l]{--},
  morecomment=[s]{/-}{-/},
}
\lstset{
  backgroundcolor=\color{codegray},
  basicstyle=\ttfamily\scriptsize,
  keywordstyle=\color{leankw}\bfseries,
  commentstyle=\color{leancomment}\itshape,
  breaklines=true, frame=single, rulecolor=\color{gray!50},
  language=lean, numbers=left, numberstyle=\tiny\color{gray},
  xleftmargin=1.5em, framexleftmargin=1.5em,
}

\newcommand{\cbrt}[1]{\sqrt[3]{#1}}
\newcommand{\prt}[2]{\sqrt[#1]{#2}}
\newcommand{\eps}{\varepsilon}

\title{\textbf{A Formally Verified Pandrosion--Steffensen Scheme}\\[0.3em]
  \large Closed-form real-line analysis and complex multi-start extension\\
  for the equation $z^p = x$, with a Lean 4 / Mathlib v4.7.0 development}
\author{Ivan \textsc{Besevic}\\
\small Independent Mathematical Research\\
\small \texttt{github.com/ivan-fr/poussiere\_de\_x}}
\date{April 2026}

\begin{document}

\maketitle

\begin{abstract}
This paper combines, in a single document, the real-line analysis of
Pandrosion's classical geometric construction for $p$-th roots and its
formally verified complex-multi-start extension.

\medskip

\textbf{Part I (real line).} We revisit Pandrosion's construction
reported by Pappus for approximating $\cbrt{2}$ and generalize it to
the computation of $p$-th roots of any positive real $x$ by inserting
proportional means via Thales' theorem. After normalization the
iteration reduces to the universal form
$s_{n+1} = 1 - (x-1)/(x\,S_p(s_n))$ with
$S_p(s) = 1 + s + \dots + s^{p-1}$, fixed point $s^* = x^{-1/p}$, and
readout $v_n = x\,s_n^{p-1} \to x^{1/p}$. We derive the closed form
\[
  \lambda_{p,x}
  = \frac{(\alpha - 1)\sum_{k=1}^{p-1} k\,\alpha^{p-1-k}}
         {\sum_{k=0}^{p-1} \alpha^k},
  \qquad \alpha = x^{1/p},
\]
for the linear contraction ratio, prove it is monotone in both $p$
and $x$ uniformly (via the closed-form identities $P_p(\alpha) = p\,D_p'(\alpha)$ and $B_p(\alpha) = \sum_{j=1}^{p} j\alpha^{j-1}$, both with all positive coefficients), with explicit large-$p$ limit $1 - \ln(x)/(x-1)$, give
non-asymptotic error bounds, extend the theory to $x < 1$ with an
explicit threshold $\alpha^*_p$ below which the iteration diverges,
and add two accelerations: a canonical starting point
$s_0 = h(1) = 1 - (x-1)/(xp)$, and Aitken--Steffensen $\Delta^2$
yielding quadratic convergence with a much smaller constant than
Newton. A scaling preconditioning $x = A \cdot x'$ exploits the
monotonicity of $\lambda$ to collapse iteration counts for large $x$.

\medskip

\textbf{Part II (complex extension, formally verified).} We extend the
iterator $h_{p,x}$ to~$\mathbb{C}$, where it has $p$ attractive fixed
points --- the $p$-th roots of $1/x$ --- separated by a non-trivial
Julia structure. A multi-start scheme
$\mathrm{seed}_{\alpha,\varepsilon}(z) = \alpha + \varepsilon(z-\alpha)$
with target-dependent $\varepsilon$ places any $z \in \mathbb{C}$
inside the quadratic-contraction disk of a chosen anchor $\alpha$,
providing a universal convergence theorem with closed-form iteration
count $N_{\mathrm{eff}} \sim \log_2 \log(1/\varepsilon)$. A
bundling theorem (Theorem~\ref{thm:grand-master}, the ``grand master''
of the Lean formalization) collects target selection, Steffensen fixing,
scaling compatibility, reconstruction, and optimal seed into a single
glue result that downstream Lean modules cite by name.
Unconditionally (without seeding), each anchor's basin of attraction
contains an explicit disk of radius $R_\sigma(x, p, \gamma_k)$, giving
the lower bound $\mathrm{vol}(\bigcup_k \Omega_k) \geq \sum_k \pi R_\sigma(x, p, \gamma_k)^2$
on the total basin volume (Proposition~\ref{prop:basin-lower-bound}). The
induced solution map is Borel-measurable everywhere and continuous
outside the Voronoï boundary of the cyclotomic anchor set (which is
Lebesgue-null). All of Part II is formally verified in a Lean 4 /
Mathlib v4.7.0 development of 121 modules with zero \texttt{sorry}
and audited axiomatic closure reducing to the standard Lean whitelist.

\medskip

The paper makes no claim of progress on classical open problems of
polynomial root-finding; its main value is to record a concrete
derivative-free root-finding scheme with closed-form convergence
analysis and complete formal verification for the monomial family
$z^p = x$.
\end{abstract}

\tableofcontents
\newpage

%=====================================================================
\part{Real-line theory}
%=====================================================================

%=====================================================================
\section{Introduction}
%=====================================================================

The Greek problem of doubling the cube consists of constructing a segment of length $\cbrt{2}$ from a unit segment. It is equivalent to finding two proportional means $m_1, m_2$ between $1$ and $2$, that is,
\[
  \frac{1}{m_1} = \frac{m_1}{m_2} = \frac{m_2}{2},
\]
which yields $m_1 = \cbrt{2}$ and $m_2 = \cbrt{4}$. Pandrosion of Alexandria proposed an \emph{approximate} geometric method, reported by Pappus in the \textit{Collectio}, Book~III~\cite{pappus}: by successive parallels (Thales' theorem), she constructs proportional ratios that approach these means without ever reaching them exactly in a finite number of steps.

\medskip

This work starts from a fundamental observation:

\begin{quote}
\itshape Mathematically the target number is exact, but any approximation procedure yields a measurable residual profile --- a footprint intrinsic to the chosen method.
\end{quote}

At each finite step of Pandrosion's construction, there remains a gap $\eps_n = |v_n - \alpha|$ between the geometric approximation and the ideal value. The sequence $(\eps_n)$ is the \textbf{residual profile} of the method. Our goal is to compute this profile exactly for any $p$-th root $x^{1/p}$, and to identify the contraction ratio that characterizes Pandrosion's method.

\medskip

\paragraph{Positioning and related work.}
Rational iterations for $p$-th roots are a classical topic with a large literature: Newton's method applied to $f(u) = u^p - x$, Halley's method, and more generally Schröder's (1870) and König's (1884) families produce iterations of arbitrary order with explicit asymptotic constants. The systematic analysis of fixed-point iterations and their acceleration goes back to Ostrowski~\cite{ostrowski}, Traub~\cite{traub}, and Householder~\cite{householder}; the derivative-free counterpart via Aitken--Steffensen (1926/1933) is standard and appears in every modern numerical-analysis textbook~\cite{stoer,numrecipes}.
What is specific to the iteration~\eqref{eq:pandrosion_gen} below is that it arises directly from Pandrosion's geometric construction, \emph{without} being derived as Newton or Schröder applied to a particular polynomial: it is a rational iteration with linear (not quadratic) order, whose contraction ratio admits the closed form~\eqref{eq:lambda_general} as a rational function of $\alpha = x^{1/p}$.
We have not located this closed form, together with its monotonicity in $(p,x)$ and its use to drive a scaling preconditioning, in the textbooks~\cite{stoer,numrecipes} or the historical literature on cube duplication, but we make no claim that the iteration itself is genuinely new: the contribution is to record the profile of a specific historical iteration and to show how standard acceleration devices (Aitken $\Delta^2$, scaling) interact with it.

\medskip

\textbf{Outline.} Section~\ref{sec:pandrosion3} recalls Pandrosion's construction for $\cbrt{2}$ and analyzes its residual profile. Section~\ref{sec:generalisation} states the generalization to any $p, x$. Section~\ref{sec:coefficient} computes the contraction ratio $\lambda_{p,x}$ and provides closed-form expressions. Section~\ref{sec:exemples} presents numerical tables. Section~\ref{sec:comparaisons} compares with bisection, the secant method, and Newton's method. Section~\ref{sec:functional} establishes functional properties of $\lambda_{p,x}$. Section~\ref{sec:nonasymptotic} provides non-asymptotic error bounds. Section~\ref{sec:extension_x_less_1} extends the theory to $x < 1$. Section~\ref{sec:optimal_start} determines the optimal starting point. Section~\ref{sec:steffensen} shows that Steffensen's acceleration of Pandrosion achieves quadratic convergence. Section~\ref{sec:scaling} exploits the monotonicity of $\lambda_{p,x}$ to optimize the method for large~$x$ via scaling. Section~\ref{sec:conclusion} concludes.

%=====================================================================
\section{Pandrosion's construction for \texorpdfstring{$\cbrt{2}$}{cbrt(2)}}
\label{sec:pandrosion3}
%=====================================================================

\subsection{The two sequences}

The geometric construction produces two sequences as a function of an initial parameter $u_0$:
\begin{equation}
  \label{eq:pandrosion_orig}
  u_{n+1} = \frac{u_n}{2 - 2\!\left(\dfrac{4 - u_n}{4}\right)^{\!3}},
  \qquad
  v_n = 2\!\left(\frac{4 - u_n}{4}\right)^{\!2},
  \qquad u_0 = 0.5.
\end{equation}

The sequence $(u_n)$ is the \emph{internal parameter} of the construction (a length in the Thales diagram). The sequence $(v_n)$ is the \emph{readout}: the approximation of $\cbrt{2}$ produced at step $n$.

\begin{proposition}
$\displaystyle\lim_{n\to\infty} v_n = \cbrt{2}$.
\end{proposition}

\subsection{Normalized substitution}

Setting $s_n = (4 - u_n)/4$, so that $s_n \in (0,1)$, we obtain:
\begin{itemize}
  \item $u_n = 4(1 - s_n)$,
  \item $v_n = 2s_n^2$,
  \item $u_{n+1} = \dfrac{4(1-s_n)}{2(1-s_n^3)} = \dfrac{2}{1 + s_n + s_n^2}$.
\end{itemize}

The iteration on $s_n$ then becomes:
\begin{equation}
  \label{eq:iter_s3}
  s_{n+1} = 1 - \frac{u_{n+1}}{4} = \frac{2s_n^2 + 2s_n + 1}{2s_n^2 + 2s_n + 2}.
\end{equation}

\subsection{Fixed point}

The fixed point $s^*$ satisfies $s^* = (2s^{*2} + 2s^* + 1)/(2s^{*2} + 2s^* + 2)$, which gives $2(1-s^*)(1+s^*+s^{*2}) = 1$, i.e.\ $2(1-s^{*3}) = 1$, hence $s^{*3} = 1/2$, so:
\[
  s^* = 2^{-1/3} = \frac{1}{\cbrt{2}}.
\]
The output at the fixed point is $v^* = 2s^{*2} = 2 \cdot 2^{-2/3} = 2^{1/3} = \cbrt{2}$. \checkmark

\subsection{Residual profile: linear convergence}

Let $h(s) = (2s^2 + 2s + 1)/(2s^2 + 2s + 2)$. The derivative at $s^*$ gives the convergence rate:
\[
  h'(s^*) = \frac{4s^* + 2}{(2s^{*2} + 2s^* + 2)^2} = \frac{2(2s^*+1)}{4(s^{*2}+s^*+1)^2} = \frac{2s^*+1}{2(s^{*2}+s^*+1)^2}.
\]

Substituting $s^* = 2^{-1/3}$ and using the closed form~\eqref{eq:lambda3}:
\[
  \lambda_{3,2} = \frac{(\cbrt{2}-1)(\cbrt{2}+2)}{\cbrt{4}+\cbrt{2}+1} \approx 0.2202.
\]

\begin{theorem}[Linear residual law for $p=3$, $x=2$]
Under the above hypotheses, for any $s_0$ sufficiently close to $s^*$:
\[
  \eps_{n+1}^{(s)} \approx \lambda_{3,2}\,\eps_n^{(s)}, \qquad \lambda_{3,2} \approx 0.2202.
\]
For the output, a Taylor expansion of $v(s) = x s^{p-1}$ at $s^* = \alpha^{-1}$ gives $v_n - \alpha \approx x(p-1)(s^*)^{p-2}(s_n - s^*)$. Using $x(s^*)^{p-2} = \alpha^p \cdot \alpha^{-(p-2)} = \alpha^2$, this rewrites as $v_n - \alpha \approx (p-1)\alpha^2 (s_n - s^*)$; the two forms are equal. For $(p,x) = (3,2)$ the prefactor is $2 \cdot 2^{2/3} = 2\,\cbrt{4}$, and the output residual of $v_n$ converges at the same linear rate $\lambda_{3,2}$.
\end{theorem}

\begin{table}[h]
\centering
\begin{tabular}{ccccc}
\toprule
$n$ & $s_n$ & $v_n$ & $\eps_n = |v_n - \cbrt{2}|$ & $\eps_{n+1}/\eps_n$ \\
\midrule
0 & 0.8750 & 1.5312 & $2.713 \times 10^{-1}$ & --- \\
1 & 0.8107 & 1.3143 & $5.44\phantom{0} \times 10^{-2}$ & 0.200 \\
2 & 0.7974 & 1.2717 & $1.17\phantom{0} \times 10^{-2}$ & 0.216 \\
3 & 0.7945 & 1.2625 & $2.58\phantom{0} \times 10^{-3}$ & 0.219 \\
4 & 0.7939 & 1.2605 & $5.67\phantom{0} \times 10^{-4}$ & 0.220 \\
$\infty$ & $2^{-1/3}$ & $\cbrt{2}$ & $0$ & $\to \lambda_{3,2}$ \\
\bottomrule
\end{tabular}
\caption{Numerical residual profile for $p=3$, $x=2$ ($u_0 = 0.5$, i.e.\ $s_0 = 0.875$). The ratio converges to $\lambda_{3,2} \approx 0.220$.}
\end{table}

%=====================================================================
\section{Geometric origin: why $S_p$ arises naturally}
\label{sec:geometrie}
%=====================================================================

\subsection{The construction by successive parallels}

Pandrosion's construction relies on Thales' theorem applied to a rectangle. Consider a rectangle with sides $H$ and $W$ equipped with two diagonals (Figure~\ref{fig:geometry}). A horizontal line at height $u_n$ from the bottom intersects the two diagonals and defines:
\begin{itemize}
  \item the parameter $s_n = (H - u_n)/H$ (normalized position from the top),
  \item the readout length $v_n = x \cdot s_n^{p-1}$ (the output of the construction).
\end{itemize}

\begin{figure}[H]
  \centering
  \IfFileExists{pandrosion_geometry.pdf}
    {\includegraphics[width=0.52\textwidth]{pandrosion_geometry.pdf}}
    {\fbox{\parbox{0.52\textwidth}{\centering \vspace{2em}\textit{Figure pandrosion\_geometry.pdf not found.}\vspace{2em}}}}
  \caption{Pandrosion's geometric construction for $p=3$, $x=2$.
  The blue horizontal lines (steps $n=0,1,2,3$) converge toward the dashed red line,
  which corresponds to $v^* = \sqrt[3]{2}$.}
  \label{fig:geometry}
\end{figure}

\subsection{Where does the polynomial $S_p$ come from?}

The fundamental idea is as follows. We seek to insert $p-1$ proportional means between $1$ and $x$, that is, to construct the ideal geometric sequence $1, \alpha, \alpha^2, \ldots, \alpha^{p-1}, x$ where $\alpha = x^{1/p}$. If $s_n \approx \alpha^{-1}$ is the current approximation of the inverse ratio, the sum of the approximate geometric sequence is:
\[
  S_p(s_n) = 1 + s_n + s_n^2 + \cdots + s_n^{p-1}.
\]
The ideal sequence satisfies $x \cdot S_p(s^*) \cdot (1-s^*) = x-1$, which rewrites as $x(1-s^{*p}) = x-1$. Pandrosion's correction adjusts $s_n$ so as to reduce the gap between $S_p(s_n)$ and its ideal value via the formula:
\[
  s_{n+1} = 1 - \frac{x-1}{x \cdot S_p(s_n)}.
\]
Geometrically, the ratio $(x-1)/(x \cdot S_p(s_n))$ measures the \emph{deviation} of the current proportional chain from the perfect geometric chain, and subtraction from $1$ produces the corrected ratio. The polynomial $S_p$ is thus the analytical transcription of the sum of the $p$ segments constructed via Thales' theorem.

%=====================================================================
\section{Generalization: Pandrosion's construction for \texorpdfstring{$x^{1/p}$}{x\^{}(1/p)}}
\label{sec:generalisation}
%=====================================================================

\subsection{The universal iteration}

Let $x > 0$ and $p \geq 2$ be an integer. Set $\alpha = x^{1/p}$. The generalized Pandrosion geometric construction is defined as follows.

\begin{definition}[Generalized Pandrosion iteration]
\label{def:pandrosion_gen}
We define the sequence $(s_n)_{n \geq 0}$ in $(0, 1)$ by
\begin{equation}
  \label{eq:pandrosion_gen}
  s_{n+1} = 1 - \frac{x - 1}{x \cdot S_p(s_n)},
  \qquad
  S_p(s) = \sum_{k=0}^{p-1} s^k = \frac{1 - s^p}{1 - s},
\end{equation}
and the output sequence by
\begin{equation}
  \label{eq:vn_gen}
  v_n = x \cdot s_n^{p-1}.
\end{equation}
We call $\eps_n = |v_n - \alpha|$ the \textbf{Pandrosion residual} at step $n$.
\end{definition}

\begin{remark}
For $p = 3$, $x = 2$, we recover exactly the iteration~\eqref{eq:iter_s3}. Indeed:
\[
  1 - \frac{1}{2(1+s+s^2)} = \frac{2s^2+2s+1}{2s^2+2s+2}. \checkmark
\]
\end{remark}

\subsection{Geometric interpretation}

The iteration~\eqref{eq:pandrosion_gen} analytically translates the Thales operation: at each step, $S_p(s_n)$ measures the ``geometric sum'' of the constructed segments, and the correction $(x-1)/(x \cdot S_p(s_n))$ adjusts the position proportionally to the gap with the ideal geometric sequence $1, \alpha, \alpha^2, \ldots, \alpha^{p-1}, x$.

\subsection{Fixed point}

\begin{theorem}[Fixed point]
\label{thm:fixed_point}
The unique fixed point of the iteration~\eqref{eq:pandrosion_gen} in $(0,1)$ is $s^* = x^{-1/p} = \alpha^{-1}$. The corresponding output is $v^* = x \cdot (x^{-1/p})^{p-1} = x^{1/p} = \alpha$.
\end{theorem}

\begin{proof}
At a fixed point $s^*$, we have $1 - (x-1)/(x \cdot S_p(s^*)) = s^*$, i.e.\
\[
  x \cdot S_p(s^*) \cdot (1 - s^*) = x - 1.
\]
Since $S_p(s)(1-s) = 1 - s^p$, this gives $x(1 - s^{*p}) = x - 1$, hence $s^{*p} = 1/x$, i.e.\ $s^* = x^{-1/p}$. The formula for $v^*$ follows immediately.
\end{proof}

%=====================================================================
\section{The contraction ratio \texorpdfstring{$\lambda_{p,x}$}{lambda}}
\label{sec:coefficient}
%=====================================================================

\subsection{General computation}

\begin{theorem}[Pandrosion contraction ratio]
\label{thm:lambda}
The convergence of the iteration~\eqref{eq:pandrosion_gen} toward $s^* = \alpha^{-1}$ is \emph{linear} with contraction ratio
\begin{equation}
  \label{eq:lambda_general}
  \boxed{
  \lambda_{p,x} = \frac{(\alpha - 1)\displaystyle\sum_{k=1}^{p-1} k\,\alpha^{p-1-k}}{\displaystyle\sum_{k=0}^{p-1} \alpha^k}
  }, \qquad \alpha = x^{1/p}.
\end{equation}
In particular, the residuals obey the law $\eps_{n+1} \approx \lambda_{p,x} \cdot \eps_n$.
\end{theorem}

\begin{proof}
Let $h(s) = 1 - (x-1)/(x \cdot S_p(s))$. Then
\[
  h'(s) = \frac{x-1}{x} \cdot \frac{S_p'(s)}{S_p(s)^2},
\]
where $S_p'(s) = \sum_{k=1}^{p-1} k s^{k-1}$.

At the fixed point $s^* = \alpha^{-1}$:
\begin{align*}
  S_p(s^*) &= \frac{1 - \alpha^{-p}}{1 - \alpha^{-1}} = \frac{\alpha(x-1)}{x(\alpha-1)}, \\
  S_p'(s^*) &= \sum_{k=1}^{p-1} k\,\alpha^{-(k-1)} = \alpha \sum_{k=1}^{p-1} k\,\alpha^{-k}.
\end{align*}

Substituting:
\[
  h'(s^*) = \frac{x-1}{x} \cdot \frac{\alpha \sum_{k=1}^{p-1} k\,\alpha^{-k}}{\left[\dfrac{\alpha(x-1)}{x(\alpha-1)}\right]^2}
  = \frac{x(\alpha-1)^2}{\alpha(x-1)} \sum_{k=1}^{p-1} k\,\alpha^{-k}.
\]

Multiplying numerator and denominator by $\alpha^{p-1}$, and using $x = \alpha^p$:
\[
  h'(s^*) = \frac{(\alpha-1)^2\,\alpha^{p-1}\sum_{k=1}^{p-1}k\,\alpha^{-k}}{\alpha^p - 1}
  = \frac{(\alpha-1)\sum_{k=1}^{p-1}k\,\alpha^{p-1-k}}{\sum_{k=0}^{p-1}\alpha^k},
\]
where we used $\alpha^p - 1 = (\alpha-1)\sum_{k=0}^{p-1}\alpha^k$. This establishes~\eqref{eq:lambda_general}.
\end{proof}

\begin{remark}
The contraction ratio $\lambda_{p,x} = h'(s^*)$ is entirely determined by $\alpha = x^{1/p}$: the target itself dictates the rate at which Pandrosion's method approaches it.
\end{remark}

\subsection{Closed forms for \texorpdfstring{$p=2$}{p=2} and \texorpdfstring{$p=3$}{p=3}}

For $p = 2$: $\sum_{k=1}^{1} k\,\alpha^{0} = 1$, $\sum_{k=0}^{1}\alpha^k = 1 + \alpha$.
\begin{equation}
  \label{eq:lambda2}
  \lambda_{2,x} = \frac{\alpha - 1}{\alpha + 1} = \frac{\sqrt{x} - 1}{\sqrt{x} + 1}.
\end{equation}

For $p = 3$: $\sum_{k=1}^{2} k\,\alpha^{2-k} = \alpha + 2$, $\sum_{k=0}^{2}\alpha^k = \alpha^2+\alpha+1$.
\begin{equation}
  \label{eq:lambda3}
  \lambda_{3,x} = \frac{(\alpha-1)(\alpha+2)}{\alpha^2 + \alpha + 1} = \frac{(\cbrt{x}-1)(\cbrt{x}+2)}{\cbrt{x^2}+\cbrt{x}+1}.
\end{equation}

For $p = 4$: $\sum_{k=1}^{3}k\,\alpha^{3-k} = \alpha^2 + 2\alpha + 3$, $\sum_{k=0}^{3}\alpha^k = \alpha^3+\alpha^2+\alpha+1$.
\begin{equation}
  \label{eq:lambda4}
  \lambda_{4,x} = \frac{(\alpha-1)(\alpha^2+2\alpha+3)}{\alpha^3+\alpha^2+\alpha+1}.
\end{equation}

\begin{remark}[Asymptotic behavior in $x$]
For any fixed $p \geq 2$ and $x \to \infty$, $\lambda_{p,x} \to 1^-$ (the dominant terms in numerator and denominator of~\eqref{eq:lambda_general} are both of order $\alpha^{p-1}$ with a ratio tending to~$1$). The larger $x$, the slower the convergence. For $x \to 1^+$, $\lambda_{p,x} \to 0$ by Proposition~\ref{prop:limits}(1): convergence is very fast near $1$.
\end{remark}

\subsection{Basin of attraction}

\begin{theorem}[Global convergence]
\label{thm:bassin}
For any $x > 1$ and any $p \geq 2$, the iteration~\eqref{eq:pandrosion_gen} converges to $s^*$ for all $s_0 \in (0,1)$. Moreover, convergence is \emph{monotone}: if $s_0 < s^*$ then $(s_n)$ is strictly increasing, and if $s_0 > s^*$ then $(s_n)$ is strictly decreasing.
\end{theorem}

\begin{proof}
Set $h(s) = 1 - (x-1)/(x \cdot S_p(s))$.

\textbf{Step 1: $h$ is strictly increasing.}
We have $h'(s) = \frac{x-1}{x} \cdot \frac{S_p'(s)}{S_p(s)^2}$, where $S_p'(s) = \sum_{k=1}^{p-1} k s^{k-1} > 0$ for $s > 0$.

\textbf{Step 2: $h$ maps $(0,1)$ into itself.}
For $s \in (0,1)$, $S_p(s) > 1$ since the first term equals $1$ and the remaining terms are strictly positive. Hence
$(x-1)/(x\,S_p(s)) < (x-1)/x < 1$, which gives $h(s) > 0$. On the other hand $(x-1)/(x\,S_p(s)) > 0$ for $x > 1$, so $h(s) < 1$. Thus $h((0,1)) \subset (0,1)$. (At the boundary: $h(0^+) = 1 - (x-1)/x = 1/x \in (0,1)$ and $h(1^-) = 1 - (x-1)/(xp) \in (0,1)$.)

\textbf{Step 3: uniqueness of the fixed point on $(0,1)$.}
The fixed-point equation $h(s) = s$ rewrites as $x\,S_p(s)(1-s) = x - 1$, i.e.\ $x(1 - s^p) = x - 1$, i.e.\ $s^p = 1/x$. On $(0,1)$ this has the unique solution $s^* = x^{-1/p}$. In particular, the sign of $h(s) - s$ is constant on each of $(0, s^*)$ and $(s^*, 1)$.

\textbf{Step 4: the sign of $h(s) - s$.}
At $s = 1/x \in (0, s^*)$ (using $x > 1$ and $s^* = x^{-1/p} > 1/x$), a one-line computation:
\[
  h(1/x) - 1/x
  \;=\; 1 - \frac{x-1}{x\, S_p(1/x)} - \frac{1}{x}
  \;=\; \frac{x-1}{x} - \frac{x-1}{x\, S_p(1/x)}
  \;=\; \frac{x-1}{x} \cdot \frac{S_p(1/x) - 1}{S_p(1/x)}
  \;>\; 0,
\]
since $S_p(1/x) > 1$ for $x > 1$ (the first term equals 1, the remaining terms are strictly positive). Hence $h(s) > s$ on $(0, s^*)$ by Step~3 (sign is constant on each sub-interval). By symmetry, $h(s) < s$ on $(s^*, 1)$.

\textbf{Step 5: monotone convergence.}
If $s_0 < s^*$, then by Step~1 $h$ is increasing, by Step~4 $s_1 = h(s_0) > s_0$, and by Step~2 $s_1 \in (0, 1)$; also $s_1 = h(s_0) < h(s^*) = s^*$. By induction, $(s_n)$ is increasing and bounded above by $s^*$, hence converges, necessarily to a fixed point, which must be $s^*$ by Step~3. The case $s_0 > s^*$ is symmetric.
\end{proof}

\subsection{Bound on the number of iterations}

\begin{proposition}[Cost of Pandrosion's method]
\label{prop:cout}
To achieve accuracy $|v_n - \alpha| < \delta$ starting from $s_0 \in (0,1)$, the required number of iterations satisfies
\[
  n_\delta \leq \left\lceil \frac{\ln(\eps_0 / \delta)}{\ln(1/\lambda_{p,x})} \right\rceil,
\]
where $\eps_0 = |v_0 - \alpha|$ is the initial error.
\end{proposition}

\begin{proof}
Since $\eps_{n+1} \leq \lambda_{p,x}\,\eps_n$ asymptotically, we have $\eps_n \leq \lambda_{p,x}^n \eps_0$. The inequality $\lambda_{p,x}^n \eps_0 < \delta$ gives $n > \ln(\eps_0/\delta)/\ln(1/\lambda_{p,x})$.
\end{proof}

\begin{example}
For $p=3$, $x=2$ ($\lambda_{3,2} \approx 0.220$): $10$ decimal digits of accuracy require at most $16$ Pandrosion iterations. For $p=3$, $x=8$ ($\lambda_{3,8} = 4/7$): $43$ iterations are needed.
\end{example}

%=====================================================================
\section{Graphical illustrations}
\label{sec:figures}
%=====================================================================

\begin{figure}[H]
  \centering
  \IfFileExists{pandrosion_figures.pdf}
    {\includegraphics[width=\textwidth]{pandrosion_figures.pdf}}
    {\fbox{\parbox{0.9\textwidth}{\centering \vspace{2em}\textit{Figure pandrosion\_figures.pdf not found.}\vspace{2em}}}}
  \caption{\textbf{Left:} Cobweb diagram of the iteration $s_{n+1} = h(s_n)$ for
  $p=3$, $x=2$, starting at $s_0=0.875$. The spiral converges to $s^* = 2^{-1/3}$
  with rate $\lambda_{3,2} \approx 0.220$.
  \textbf{Right:} Residual profiles on a semi-logarithmic scale.
  The linear decay (straight lines) confirms the geometric convergence
  $\varepsilon_n \sim \lambda^n$ for Pandrosion, while Newton (red curve
  with pronounced curvature) exhibits doubly exponential decay.}
  \label{fig:figures}
\end{figure}

%=====================================================================
\section{Numerical examples}
\label{sec:exemples}
%=====================================================================

\subsection{Table of contraction ratios}

\begin{table}[h]
\centering
\begin{tabular}{ccccc}
\toprule
$p$ & $x$ & $\alpha = x^{1/p}$ & $\lambda_{p,x}$ (exact) & $\lambda_{p,x}$ (numerical) \\
\midrule
2 & 2 & $\sqrt{2}$ & $(\sqrt{2}-1)/(\sqrt{2}+1)$ & $0.1716$ \\
2 & 3 & $\sqrt{3}$ & $(\sqrt{3}-1)/(\sqrt{3}+1)$ & $0.2679$ \\
2 & 5 & $\sqrt{5}$ & $(\sqrt{5}-1)/(\sqrt{5}+1)$ & $0.3820$ \\
\midrule
3 & 2 & $\cbrt{2}$ & $(\cbrt{2}-1)(\cbrt{2}+2)/(\cbrt{4}+\cbrt{2}+1)$ & $0.2202$ \\
3 & 3 & $\cbrt{3}$ & $(\cbrt{3}-1)(\cbrt{3}+2)/(\cbrt{9}+\cbrt{3}+1)$ & $0.3366$ \\
3 & 8 & $2$ & $(2-1)(2+2)/(4+2+1)$ & $0.5714$ \\
\midrule
4 & 2 & $\sqrt[4]{2}$ & (formula~\eqref{eq:lambda4}) & $0.2427$ \\
4 & 16 & $2$ & $(2-1)(4+4+3)/(8+4+2+1)$ & $0.7333$ \\
\midrule
5 & 2 & $\sqrt[5]{2}$ & (general formula) & $0.2565$ \\
\bottomrule
\end{tabular}
\caption{Contraction ratios $\lambda_{p,x}$ for various values of $p$ and $x$. A value close to $0$ means very fast convergence; close to $1$, very slow.}
\end{table}

\subsection{Detailed numerical profile for \texorpdfstring{$p=2$, $x=2$}{p=2, x=2}}

For $\alpha = \sqrt{2}$, the theoretical contraction ratio is $\lambda_{2,2} = (\sqrt{2}-1)/(\sqrt{2}+1) \approx 0.1716$ (formula~\eqref{eq:lambda2}). The iteration is:
\[
  s_{n+1} = 1 - \frac{1}{2(1 + s_n)}, \qquad v_n = 2s_n.
\]

\begin{table}[h]
\centering
\begin{tabular}{ccccc}
\toprule
$n$ & $s_n$ & $v_n$ & $\eps_n = |v_n - \sqrt{2}|$ & $\eps_{n+1}/\eps_n$ \\
\midrule
0 & 0.7500 & 1.5000 & $8.579\times 10^{-2}$ & --- \\
1 & 0.7143 & 1.4286 & $1.44\phantom{0}\times 10^{-2}$ & 0.167 \\
2 & 0.7083 & 1.4167 & $2.45\phantom{0}\times 10^{-3}$ & 0.171 \\
3 & 0.7073 & 1.4146 & $4.21\phantom{0}\times 10^{-4}$ & 0.171 \\
4 & 0.7071 & 1.4143 & $7.22\phantom{0}\times 10^{-5}$ & 0.172 \\
5 & 0.7071 & 1.4142 & $1.24\phantom{0}\times 10^{-5}$ & 0.171 \\
$\infty$ & $1/\sqrt{2}$ & $\sqrt{2}$ & $0$ & $\to \lambda_{2,2}$ \\
\bottomrule
\end{tabular}
\caption{Numerical profile for $p=2$, $x=2$, $s_0 = 0.75$. The ratio $\eps_{n+1}/\eps_n$ converges rapidly to $\lambda_{2,2} \approx 0.172$.}
\end{table}

\subsection{The residual as a characterization of a target}

Formula~\eqref{eq:lambda_general} implies a useful observation: $\lambda_{p, \alpha^p} = f_p(\alpha)$ where $f_p$ is a rational function of $\alpha$ alone. In particular, the residual profile depends only on the value of the sought root $\alpha$ and not on its representation as $x = \alpha^p$.

In particular, for $p = 3$:
\[
  \lambda_{3,x} = \frac{(\alpha-1)(\alpha+2)}{\alpha^2+\alpha+1} \xrightarrow[\alpha \to 1]{} 0, \quad
  \lambda_{3,x} \xrightarrow[\alpha \to \infty]{} 1.
\]

%=====================================================================
\section{Comparison with other methods}
\label{sec:comparaisons}
%=====================================================================

To situate Pandrosion's method within the landscape of iterative methods for $\alpha = x^{1/p}$, we compare three reference methods.

\subsection{Newton's method}

Newton's method for $f(u) = u^p - x$ gives the iteration
\[
  u_{n+1} = \frac{(p-1)u_n + x/u_n^{p-1}}{p},
\]
whose convergence to $\alpha = x^{1/p}$ is \emph{quadratic}:
\[
  \eps_{n+1}^{\text{Newton}} \approx K_{p,x}\,(\eps_n^{\text{Newton}})^2, \qquad K_{p,x} = \frac{p-1}{2\alpha}.
\]

\subsection{Bisection}

Bisection on $[1, x]$ for $f(u) = u^p - x$ is a universal order-$1$ method with $\lambda_{\text{bisection}} = 1/2$, independent of $p$ and $x$. Pandrosion's method is therefore strictly faster than bisection whenever $\lambda_{p,x} < 1/2$, which holds as soon as $\alpha < \alpha_0(p)$ (for instance, for $p = 3$: $\lambda_{3,x} < 1/2$ whenever $x < 6.75$).

\subsection{Secant method}

The secant method for $f(u) = u^p - x$ has order $\varphi = (1+\sqrt{5})/2 \approx 1.618$ (superlinear). For $p = 3$, $x = 2$, it reaches machine precision in about $6$ steps, compared to $\sim 16$ for Pandrosion and $\sim 5$ for Newton.

\subsection{Comparative table}

\begin{table}[h]
\centering
\renewcommand{\arraystretch}{1.4}
\begin{tabular}{lcccc}
\toprule
& \textbf{Pandrosion} & \textbf{Bisection} & \textbf{Secant} & \textbf{Newton} \\
\midrule
Nature & Geometric & Universal & Analytic & Analytic \\
Order $q$ & $1$ & $1$ & $\varphi \approx 1.62$ & $2$ \\
Rate ($p\!=\!3, x\!=\!2$) & $\lambda \approx 0.220$ & $\lambda = 0.500$ & superlinear & quadratic \\
Iter.\ for $10$ digits & $\sim 16$ & $\sim 34$ & $\sim 6$ & $\sim 5$ \\
Derivative required & No & No & No & Yes \\
\bottomrule
\end{tabular}
\caption{Comparison of approximation methods for $x^{1/p}$ with $p=3$, $x=2$.}
\end{table}

\subsection{Discussion}

Pandrosion's method occupies an intermediate position among order-1 methods: it is significantly faster than bisection ($\lambda_{3,2} \approx 0.22$ versus $0.50$) while remaining derivative-free. Compared to Newton and its variants, however, plain Pandrosion is slower: its order is $1$, whereas Newton attains $2$, and even with Aitken--Steffensen acceleration (Section~\ref{sec:steffensen}) the practical advantage over Newton is modest and regime-dependent (see Remark~\ref{rem:fair-comparison}). We do not claim a general algorithmic advantage over Newton or its modern variants; our interest in Pandrosion is that it is the historical construction attributed to Pandrosion of Alexandria, and that its convergence profile admits the clean closed form $\lambda_{p,x}$ of Theorem~\ref{thm:lambda}.

\begin{remark}[On fair comparison]
\label{rem:fair-comparison}
The iteration counts above are per-iterate, \emph{not} per-function-evaluation. Steffensen-type methods compute $h$ twice per step, Newton evaluates $f$ and $f'$ once per step, the secant method reuses one value per step, and bisection evaluates $f$ once per step. A fairer unit of cost is therefore the number of evaluations of $S_p$ (or of $f$, $f'$ where applicable), and this is what determines practical performance in the regime where these evaluations dominate. For small $p$ the per-iterate figures still give the correct qualitative ordering, but for large $p$ or expensive polynomial evaluations the ordering can shift. We report per-iterate counts because they match the structure of the theorems (Theorem~\ref{thm:lambda} is a statement about $\lambda_{p,x}^n$, not about wall time).
\end{remark}

The coexistence of four distinct error laws for the same target $\alpha = \cbrt{2}$ illustrates a structural point: the rate at which a fixed-point iteration approaches a given number depends on the iteration, not on the number; conversely, once the iteration is fixed, the target determines the rate through $\lambda_{p,x}$.

%=====================================================================
\section{Functional properties of \texorpdfstring{$\lambda_{p,x}$}{lambda}}
\label{sec:functional}
%=====================================================================

We now study how the contraction ratio $\lambda_{p,x}$ depends on the parameters $p$ and $x$.

\subsection{Monotonicity}

\begin{theorem}[Double monotonicity]
\label{thm:monotonicity}
The contraction ratio $\lambda_{p,x}$, viewed as a function of $\alpha = x^{1/p} > 1$, satisfies:
\begin{enumerate}
  \item For fixed $p \geq 2$, $\alpha \mapsto \lambda_{p,\alpha^p}$ is strictly increasing on $(1,\infty)$.
  \item For fixed $\alpha > 1$, $p \mapsto \lambda_{p,\alpha^p}$ is strictly increasing on $\{2,3,4,\ldots\}$.
\end{enumerate}
\end{theorem}

\begin{proof}
Write $\lambda_{p,\alpha^p} = (\alpha-1) N_p(\alpha)/D_p(\alpha)$ with $N_p(\alpha) := \sum_{k=1}^{p-1} k\,\alpha^{p-1-k}$ and $D_p(\alpha) := \sum_{k=0}^{p-1}\alpha^k$. The key auxiliary fact is the following Taylor bound.

\medskip

\noindent\emph{Lemma (Taylor bound).} Let $x > 1$ and set $y := -\ln(x)/p$, $\beta := e^y = x^{-1/p}$. For $p \geq \lceil \ln x \rceil + 1$ (so that $|y| \leq 1$),
\[
  \bigl|\, p(\beta - 1) + \beta \ln x \,\bigr|
  \;\leq\; \frac{3\,(\ln x)^2}{p}.
\]
\emph{Proof of the lemma.} Since $\beta = e^y$ and $py = -\ln x$,
$p(\beta - 1) + \beta\ln x = p(\beta - 1) - p y \beta = p(\beta - 1 - y) - p y(\beta - 1)$.
For $|y| \leq 1$ we use two standard bounds:
\begin{align*}
  |e^y - 1 - y|
   &\;=\; \Bigl|\sum_{k=2}^{\infty} \frac{y^k}{k!}\Bigr|
    \;\leq\; \sum_{k=2}^{\infty} \frac{|y|^k}{k!}
    \;=\; y^2 \sum_{k=2}^{\infty} \frac{|y|^{k-2}}{k!}
    \;\leq\; y^2 \sum_{k=2}^{\infty} \frac{1}{k!}
    \;=\; (e - 2)\, y^2
    \;<\; y^2,
  \\[0.3em]
  |e^y - 1|
   &\;=\; \Bigl|\sum_{k=1}^{\infty} \frac{y^k}{k!}\Bigr|
    \;\leq\; |y| \sum_{k=1}^{\infty} \frac{|y|^{k-1}}{k!}
    \;\leq\; |y| \sum_{k=1}^{\infty} \frac{1}{k!}
    \;=\; (e - 1)\, |y|
    \;<\; 2|y|.
\end{align*}
(Here $e - 2 \approx 0.718$ and $e - 1 \approx 1.718$, so both majorants $y^2$ and $2|y|$ are loose; the loose forms suffice.) The triangle inequality then gives
\[
  |p(\beta - 1 - y) - p y(\beta - 1)|
  \;\leq\; p\, y^2 + p|y|\cdot 2|y|
  \;=\; 3 p y^2
  \;=\; \frac{3(\ln x)^2}{p}. \tag*{$\square$}
\]
(The symbol $\beta$ is the fixed-point reciprocal of the root $\alpha = x^{1/p}$ of the theorem: $\beta = 1/\alpha$.)

\medskip

\noindent Setting $\lambda_\infty(x) := 1 - \ln(x)/(x-1)$, a direct algebraic rearrangement of the closed form~\eqref{eq:lambda_general} combined with the lemma yields
\begin{equation}\label{eq:taylor-bound}
  \bigl|\, \lambda_{p,x} - \lambda_\infty(x) \,\bigr|
  \;\leq\; \frac{3\,x\,(\ln x)^2}{(x-1)\,p},
  \qquad \text{for every } p \geq \lceil \ln x \rceil + 1.
\end{equation}
(The factor $x/(x-1)$ arises as follows. For $p \geq 1$ and $x > 1$, $\ln(x)/p \leq \ln x$, so $\beta = \exp(-\ln x/p) \geq \exp(-\ln x) = 1/x$. Multiplying by $x - 1 > 0$ gives $\beta(x - 1) \geq (x-1)/x$, i.e.\ $1/[\beta(x-1)] \leq x/(x-1)$, which is the factor that appears when dividing the Taylor-lemma bound by the denominator $\beta(x-1)/(\ln x)$ of the explicit form of $\lambda_{p,x}$.)

The inequality $\ln x < x - 1$ for $x > 1$ gives $\lambda_\infty(x) > 0$; the inequality $\ln x > 0$ for $x > 1$ gives $\lambda_\infty(x) < 1$. Hence $\lambda_\infty(x) \in (0, 1)$ for every $x > 1$.

\medskip

\medskip

\noindent\emph{Key identity.} By direct telescoping (or induction on $p$),
\begin{equation}\label{eq:key-identity}
  (\alpha - 1)\, N_p(\alpha) \;=\; D_p(\alpha) - p,
\end{equation}
which is verified at small $p$: $p=2$ gives $(\alpha-1) = (1+\alpha) - 2$; $p=3$ gives $(\alpha-1)(\alpha+2) = \alpha^2 + \alpha - 2 = (1+\alpha+\alpha^2) - 3$. Differentiating~\eqref{eq:key-identity}: $N_p + (\alpha-1) N_p' = D_p'$, whence $(\alpha-1) N_p' = D_p' - N_p$.

\medskip

\noindent\emph{Part (1) --- monotonicity in $\alpha$ for fixed $p$, uniform.} After multiplying through by the positive denominator $(\alpha-1) N_p(\alpha) D_p(\alpha)$, the log-derivative
\[
  \frac{d}{d\alpha}\log\lambda_{p,\alpha^p}
  = \frac{1}{\alpha - 1} + \frac{N_p'(\alpha)}{N_p(\alpha)} - \frac{D_p'(\alpha)}{D_p(\alpha)}
\]
yields the polynomial numerator
\[
  P_p(\alpha) \;:=\; N_p D_p + (\alpha - 1)\, N_p'\, D_p - (\alpha - 1)\, N_p\, D_p'.
\]
Substituting $(\alpha-1) N_p' = D_p' - N_p$ and $(\alpha-1) N_p = D_p - p$ from the key identity:
\begin{align*}
  P_p &= N_p D_p + (D_p' - N_p) D_p - (D_p - p) D_p' \\
      &= N_p D_p + D_p' D_p - N_p D_p - D_p D_p' + p D_p'
      = p \cdot D_p'(\alpha).
\end{align*}
Since $D_p'(\alpha) = \sum_{j=1}^{p-1} j\,\alpha^{j-1} = 1 + 2\alpha + 3\alpha^2 + \cdots + (p-1)\alpha^{p-2}$ has all coefficients positive, $P_p(\alpha) > 0$ for all $\alpha > 0$ and all $p \geq 2$. Hence the log-derivative is positive on $\alpha > 1$, and $\lambda_{p,\alpha^p}$ is strictly increasing in $\alpha$. The argument is uniform: no scope restriction is needed.

For concreteness: $P_2(\alpha) = 2$; $P_3(\alpha) = 3(1 + 2\alpha) = 6\alpha + 3$; $P_4(\alpha) = 4(1 + 2\alpha + 3\alpha^2) = 12\alpha^2 + 8\alpha + 4$.

For general $p \geq 2$, we combine the Taylor bound~\eqref{eq:taylor-bound} with the monotonicity of $\lambda_\infty$. Differentiating,
\[
  \lambda_\infty'(x)
  = \frac{\ln x - (x-1)/x}{(x-1)^2}.
\]
The numerator equals $\ln x + 1/x - 1$. Its derivative is $1/x - 1/x^2 = (x-1)/x^2 > 0$ on $x > 1$, and its value at $x = 1$ is $0$; hence it is strictly positive on $(1, \infty)$, and $\lambda_\infty$ is strictly increasing there. \begin{remark}[Monotonicity of the asymptote]
The asymptote $\lambda_\infty(x) = 1 - \ln(x)/(x-1)$ is itself strictly increasing on $(1, \infty)$: $\lambda_\infty'(x) = (\ln x - (x-1)/x)/(x-1)^2 > 0$ since $\ln x + 1/x - 1 > 0$ on $x > 1$ (value $0$ at $x = 1$, derivative $(x-1)/x^2 > 0$). This is a consistency check only; the finite-$p$ monotonicity is already established above via $P_p = p\,D_p'$, without appealing to the Taylor bound.
\end{remark}

\medskip

\medskip

\noindent\emph{Part (2) --- monotonicity in $p$ for fixed $\alpha > 1$, uniform.} From the recurrences $N_{p+1} = \alpha N_p + p$ and $D_{p+1} = D_p + \alpha^p$, direct algebra gives
\[
  \lambda_{p+1,\alpha^{p+1}} - \lambda_{p,\alpha^p}
  = \frac{(\alpha-1)\cdot B_p(\alpha)}{D_p(\alpha)\,D_{p+1}(\alpha)},
  \qquad
  B_p(\alpha) := (\alpha-1)\, N_p D_p + p\, D_p - \alpha^p N_p.
\]
Using~\eqref{eq:key-identity}, $(\alpha-1) N_p D_p = (D_p - p) D_p = D_p^2 - p D_p$, so
\[
  B_p(\alpha) \;=\; D_p(\alpha)^2 \;-\; \alpha^p\, N_p(\alpha).
\]
This form has a closed-form polynomial expansion:
\begin{equation}\label{eq:Bp-closed}
  B_p(\alpha) \;=\; \sum_{j=1}^{p}\, j\, \alpha^{j-1}
  \;=\; 1 + 2\alpha + 3\alpha^2 + \cdots + p\, \alpha^{p-1},
\end{equation}
verified by induction on $p$. Base case: $B_2(\alpha) = (1+\alpha)^2 - \alpha^2 = 1 + 2\alpha$, as required. Inductive step: using the recurrence
\[
  B_{p+1}(\alpha) = \alpha\, B_p(\alpha) + D_p(\alpha) + \alpha^p
\]
(which follows from the defining recurrences of $N_{p+1}$ and $D_{p+1}$ plus the substitution $B_p = D_p^2 - \alpha^p N_p$), one computes $\alpha B_p + D_p + \alpha^p = \sum_{j=1}^{p} j\alpha^j + \sum_{k=0}^{p-1}\alpha^k + \alpha^p = \sum_{j=0}^{p} (j+1)\alpha^j = B_{p+1}$.

All coefficients of $B_p$ are strictly positive integers, so $B_p(\alpha) > 0$ for every $\alpha > 0$ and every $p \geq 2$. In particular $B_p(1) = 1 + 2 + \cdots + p = p(p+1)/2$, and $B_3(\alpha) = 3\alpha^2 + 2\alpha + 1$, $B_4(\alpha) = 4\alpha^3 + 3\alpha^2 + 2\alpha + 1$, etc. Hence $\lambda_{p+1,\alpha^{p+1}} - \lambda_{p,\alpha^p} > 0$ for all $\alpha > 1$ and $p \geq 2$, uniformly.
\end{proof}

\begin{remark}[Closed forms and numerical verification]
The closed forms $P_p(\alpha) = p \cdot D_p'(\alpha)$ and $B_p(\alpha) = \sum_{j=1}^p j\alpha^{j-1}$ emerged in the course of this write-up; a standalone script \texttt{latex/verify\_monotonicity.py} cross-checks both identities symbolically for $p \leq 20$ using SymPy, and confirms strict monotonicity of $\lambda_p(\alpha)$ for $2 \leq p \leq 20$, $\alpha \in [1.01, 100]$, and of $\lambda_p(\alpha^p)$ for $2 \leq p \leq 50$, $\alpha \in \{1.1, 1.5, 2, 3, 10, 100\}$ as a consistency check. Monotonicity itself is proven analytically above; the numerical scan is purely a consistency check.
\end{remark}

\begin{remark}
Monotonicity in $p$ has a geometric interpretation: inserting more proportional means (larger $p$) constrains the construction further, slowing convergence.
\end{remark}

\subsection{Limiting behavior}

\begin{proposition}[Asymptotic regimes]
\label{prop:limits}
\mbox{}
\begin{enumerate}
  \item \textbf{Near $\alpha = 1$} (i.e.\ $x \to 1^+$):
  \[
    \lambda_{p,x} \sim \frac{(p-1)(\alpha - 1)}{2} \quad \text{as } \alpha \to 1^+.
  \]
  \item \textbf{Large $\alpha$} (i.e.\ $x \to \infty$, $p$ fixed): $\lambda_{p,x} \to 1^-$.
  \item \textbf{Large $p$} (fixed $x > 1$): $\lambda_{p,x}$ is increasing in $p$ and
  \[
    \lim_{p \to \infty} \lambda_{p,x} \;=\; 1 - \frac{\ln x}{x - 1}
     \;\in\; (0, 1).
  \]
  In particular $\lambda_{p,x}$ does \emph{not} tend to~$1$ as $p \to \infty$ with $x$ fixed.
\end{enumerate}
\end{proposition}

\begin{proof}
For (1): $N_p$ and $D_p$ are polynomials, so $\alpha \mapsto N_p(\alpha)/D_p(\alpha)$ is continuous at $\alpha = 1$ with limit $(p-1)/2$; combined with the factor $(\alpha - 1) \to 0$, $\lambda \sim (p-1)(\alpha-1)/2$ as $\alpha \to 1^+$.

For (2): the dominant term in $N_p$ is $1 \cdot \alpha^{p-2}$ and in $D_p$ is $\alpha^{p-1}$, so $\lambda \sim (\alpha-1)/\alpha \to 1$.

For (3): the explicit limit and its convergence rate follow from the Taylor bound~\eqref{eq:taylor-bound} established in the proof of Theorem~\ref{thm:monotonicity} above:
$|\lambda_{p,x} - \lambda_\infty(x)| \leq C(x)/p$ with $C(x) = 3x(\ln x)^2/(x-1)$, valid for $p \geq \lceil \ln x \rceil + 1$. Letting $p \to \infty$ gives $\lambda_{p,x} \to \lambda_\infty(x) = 1 - \ln(x)/(x-1)$. Positivity and the upper bound $< 1$ were shown there as well.

Numerically: for $x = 2$, $\lambda_\infty(2) = 1 - \ln 2 \approx 0.3069$, and $\lambda_{200,2} \approx 0.3057$, $\lambda_{1000,2} \approx 0.3066$, $\lambda_{10000,2} \approx 0.3068$, in agreement with the bound $|\lambda_{p,2} - \lambda_\infty(2)| \leq 6(\ln 2)^2/p \approx 2.88/p$.
\end{proof}

\begin{table}[h]
\centering
\begin{tabular}{ccccc}
\toprule
$x$ & $\lambda_{2,x}$ & $\lambda_{10,x}$ & $\lambda_{100,x}$ & $1 - \ln x / (x-1)$ \\
\midrule
$1.001$ & $0.00025$ & $0.00045$ & $0.00050$ & $0.00050$ \\
$2$ & $0.1716$ & $0.2823$ & $0.3044$ & $0.3069$ \\
$10$ & $0.5195$ & $0.7123$ & $0.7412$ & $0.7442$ \\
$100$ & $0.8182$ & $0.9409$ & $0.9524$ & $0.9535$ \\
\bottomrule
\end{tabular}
\caption{Contraction ratios illustrating monotonicity in $p$ and $x$. Each column is increasing (fixed $p$); each row is increasing (fixed $x$). The rightmost column is the $p \to \infty$ limit from Proposition~\ref{prop:limits}(3); the $\lambda_{100,x}$ column is already close to this limit.}
\end{table}

%=====================================================================
\section{Non-asymptotic error bounds}
\label{sec:nonasymptotic}
%=====================================================================

Theorem~\ref{thm:lambda} gives the \emph{asymptotic} law $\eps_{n+1} \approx \lambda_{p,x} \eps_n$. We now establish a rigorous bound valid for \emph{all} $n \geq 0$.

\begin{theorem}[Non-asymptotic contraction]
\label{thm:nonasymp}
Let $x > 1$, $p \geq 2$, and $s_0 \in (0,1)$. Define
\[
  \Lambda = \Lambda(s_0, p, x) = \sup_{s \in [\min(s_0, s^*),\, \max(s_0, s^*)]} h'(s).
\]
Then $\Lambda < 1$ and for all $n \geq 0$:
\begin{equation}
  \label{eq:nonasymp_bound}
  |s_n - s^*| \leq \Lambda^n \, |s_0 - s^*|.
\end{equation}
Consequently,
\[
  |v_n - \alpha| \;\leq\; (p-1)\,\alpha^2 \cdot \Lambda^n\, |s_0 - s^*| \;+\; O(\Lambda^{2n}).
\]
The prefactor follows from a Taylor expansion of $v(s) = x\, s^{p-1}$ at $s^* = \alpha^{-1}$:
\[
  v(s) - v(s^*) \;=\; x\,(p-1)\,(s^*)^{p-2}\,(s - s^*) + O((s - s^*)^2),
\]
and $x\,(p-1)\,(s^*)^{p-2} = (p-1)\, x\, x^{-(p-2)/p} = (p-1)\, x^{2/p} = (p-1)\,\alpha^2$.
\end{theorem}

\begin{proof}
By the mean value theorem, for any $n$:
\[
  |s_{n+1} - s^*| = |h(s_n) - h(s^*)| = |h'(\xi_n)| \cdot |s_n - s^*|
\]
for some $\xi_n$ between $s_n$ and $s^*$. By Theorem~\ref{thm:bassin}, the sequence $(s_n)$ remains in $[\min(s_0,s^*), \max(s_0,s^*)]$, so $|h'(\xi_n)| \leq \Lambda$. The bound follows by induction.

That $\Lambda < 1$ follows from$h'(s^*) = \lambda_{p,x} < 1$ (Theorem~\ref{thm:lambda}) and the fact that $h'$ is continuous, bounded, and $h'(s) < 1$ for all $s \in (0,1)$ when $x > 1$ (since $h$ maps $(0,1)$ strictly into itself).
\end{proof}

\begin{remark}
The quantity $\Lambda$ can be bounded explicitly: $h'(s) = \tfrac{x-1}{x} \cdot S'_p(s)/S_p(s)^2$, so
\[
  \Lambda \;\leq\; \frac{x-1}{x} \cdot \frac{\max_{s \in I} S'_p(s)}{\min_{s \in I} S_p(s)^2},
  \qquad I := [\min(s_0, s^*), \max(s_0, s^*)],
\]
and on $(0,1)$ both $S_p$ and $S_p'$ are increasing, so the numerator is attained at the right endpoint and the denominator at the left endpoint. This yields a computable upper bound for $\Lambda$ without requiring a detailed analysis of the monotonicity of $h'$ itself.
\end{remark}

\begin{example}
For $p = 3$, $x = 2$, $s_0 = 0.5$ (so $s^* = 2^{-1/3} \approx 0.7937$), one evaluates $h'$ at both endpoints of the interval $[s_0, s^*] = [0.5, 0.7937]$:
\[
  h'(s) = \tfrac{x-1}{x} \cdot \frac{S'_3(s)}{S_3(s)^2}
        = \tfrac{1}{2} \cdot \frac{1 + 2s}{(1 + s + s^2)^2}.
\]
At $s = 0.5$: $S_3(0.5) = 1.75$, $S'_3(0.5) = 2$, so $h'(0.5) = \tfrac{1}{2} \cdot 2/(1.75)^2 \approx 0.3265$. At $s = s^*$: $h'(s^*) = \lambda_{3,2} \approx 0.2202$. Therefore
\[
  \Lambda \;=\; \max_{s \in [0.5, s^*]} h'(s) \;=\; h'(0.5) \;\approx\; 0.327,
\]
since $h'$ is decreasing on this interval. The bound~\eqref{eq:nonasymp_bound} gives $|s_n - s^*| \leq 0.327^n \cdot |0.5 - s^*| \leq 0.327^n \cdot 0.294$, so $|s_{10} - s^*| \leq 4.6 \times 10^{-6}$. The asymptotic ratio $\lambda_{3,2} \approx 0.220$ is reached after a few iterations.
\end{example}

%=====================================================================
\section{Extension to \texorpdfstring{$x < 1$}{x < 1}}
\label{sec:extension_x_less_1}
%=====================================================================

All results so far assumed $x > 1$. We now extend the theory to $0 < x < 1$.

\subsection{The fixed point leaves $(0,1)$}

For $0 < x < 1$, the target is $\alpha = x^{1/p} \in (0,1)$ and the fixed point $s^* = \alpha^{-1} > 1$. The iteration domain therefore extends beyond $(0,1)$: $s_n$ converges to $s^* > 1$.

\begin{proposition}[Local convergence for $x < 1$ with a threshold]
\label{prop:x_small}
For $0 < x < 1$ and $p \geq 2$, the formula~\eqref{eq:lambda_general} remains valid and
\[
  \lambda_{p,x}
  \;=\; \frac{(\alpha - 1)\sum_{k=1}^{p-1}k\,\alpha^{p-1-k}}{\sum_{k=0}^{p-1}\alpha^k}
  \;<\; 0
  \qquad \text{since } \alpha = x^{1/p} \in (0,1).
\]
Local convergence (i.e.\ $|\lambda_{p,x}| < 1$) does \emph{not} hold for all $x \in (0,1)$: there is a threshold $\alpha^\ast_p \in (0,1)$ such that
\[
  |\lambda_{p,x}| < 1
  \;\iff\;
  \alpha > \alpha^\ast_p.
\]
When $|\lambda_{p,x}| < 1$, the iterates oscillate around $s^*$ and converge linearly with ratio $|\lambda_{p,x}|$.
\end{proposition}

\begin{proof}
That $\lambda_{p,x} < 0$ for $\alpha < 1$ is immediate from~\eqref{eq:lambda_general}. Using $\alpha^p - 1 = (\alpha - 1) \sum_{k=0}^{p-1} \alpha^k$, one rewrites
\[
  |\lambda_{p,x}|
  \;=\; \frac{(1-\alpha)^2 \, N_p(\alpha)}{1 - \alpha^p}
  \;=\; \frac{(1 - \alpha)\, N_p(\alpha)}{D_p(\alpha)} \;=:\; F(\alpha)
\]
(using $1 - \alpha^p = (1-\alpha) D_p$), a smooth strictly positive function of $\alpha \in (0, 1)$. Boundary values: as $\alpha \to 1^-$, $F \to 0$ (the factor $1-\alpha$ goes to zero while $N_p/D_p \to (p-1)/2$ is bounded); as $\alpha \to 0^+$, $F \to 1 \cdot (p-1)/1 = p - 1 \geq 1$.

For uniqueness, we show that $F$ is strictly decreasing on $(0, 1)$. The log-derivative is
\[
  \frac{F'(\alpha)}{F(\alpha)}
  \;=\; -\frac{1}{1-\alpha} + \frac{N_p'(\alpha)}{N_p(\alpha)} - \frac{D_p'(\alpha)}{D_p(\alpha)}.
\]
The key identity~\eqref{eq:key-identity}, $(\alpha - 1) N_p = D_p - p$, and its derivative $N_p + (\alpha - 1) N_p' = D_p'$, give, for $\alpha \in (0, 1)$ (where $\alpha - 1 < 0$):
\[
  (1-\alpha)\, N_p \;=\; p - D_p,
  \qquad
  (1-\alpha)\, N_p' \;=\; N_p - D_p'.
\]
Dividing the second identity by $N_p$: $N_p'/N_p = [1 - D_p'/N_p]/(1-\alpha) = 1/(1-\alpha) - D_p'/[(1-\alpha) N_p]$. Substituting into $F'/F$:
\begin{align*}
  \frac{F'}{F}
  &= -\frac{1}{1-\alpha} + \left[\frac{1}{1-\alpha} - \frac{D_p'}{(1-\alpha) N_p}\right] - \frac{D_p'}{D_p}
   = -\frac{D_p'}{(1-\alpha) N_p} - \frac{D_p'}{D_p} \\
  &= -\frac{D_p' \cdot [D_p + (1-\alpha) N_p]}{(1-\alpha) N_p\, D_p}
   = -\frac{D_p' \cdot [D_p + p - D_p]}{(1-\alpha) N_p\, D_p}
   = -\frac{p\, D_p'(\alpha)}{(1-\alpha)\, N_p(\alpha)\, D_p(\alpha)},
\end{align*}
using $D_p + (1-\alpha) N_p = D_p + (p - D_p) = p$ in the last line. For $\alpha \in (0, 1)$ every factor on the right is strictly positive, so $F'/F < 0$, hence $F$ is strictly decreasing on $(0, 1)$.

Combined with the boundary values $F(0^+) = p - 1 \geq 1$ and $F(1^-) = 0$, $F$ crosses the level $1$ at exactly one point $\alpha^\ast_p \in (0, 1)$, and $|\lambda_{p, \alpha^p}| < 1 \iff \alpha > \alpha^\ast_p$.
\end{proof}

\begin{example}[$p = 3$ threshold]\label{ex:p3-threshold}
For $p = 3$, $|\lambda_{3,x}| = (1-\alpha)(2+\alpha)/(1 + \alpha + \alpha^2)$. Setting this equal to $1$ and simplifying yields $2\alpha^2 + 2\alpha - 1 = 0$, i.e.
\[
  \alpha^\ast_3 \;=\; \frac{-1 + \sqrt{3}}{2} \;\approx\; 0.3660,
  \qquad
  x^\ast_3 \;=\; (\alpha^\ast_3)^3 \;\approx\; 0.04904.
\]
Hence Pandrosion's iteration at $p = 3$ contracts locally iff $x > 0.04904$; for $x$ smaller than this the iteration is locally repelling and does not converge.
\end{example}

\begin{table}[h]
\centering
\begin{tabular}{cccccc}
\toprule
$p$ & $x$ & $\alpha = x^{1/p}$ & $s^* = \alpha^{-1}$ & $\lambda_{p,x}$ & $|\lambda_{p,x}|$ \\
\midrule
2 & 0.5   & $1/\sqrt{2}$   & $\sqrt{2}$    & $-0.1716$ & $0.1716$ \\
3 & 0.5   & $\cbrt{0.5}$   & $\cbrt{2}$    & $-0.2378$ & $0.2378$ \\
3 & 0.125 & $0.5$          & $2$           & $-0.7143$ & $0.7143$ \\
3 & 0.05  & $0.3684$       & $2.7144$      & $-0.9945$ & $0.9945$ \\
\midrule
3 & 0.01  & $0.2154$       & $4.6416$      & $-1.3774$ & $\mathbf{1.3774}$ \\
5 & 0.01  & $0.3981$       & $2.5119$      & $-2.0399$ & $\mathbf{2.0399}$ \\
\bottomrule
\end{tabular}
\caption{Contraction ratios for $x < 1$. The last two rows illustrate Proposition~\ref{prop:x_small}: for $x$ below the threshold $x^\ast_p$ the iteration is locally repelling ($|\lambda| > 1$) and the construction does \emph{not} converge.}
\end{table}

\begin{remark}[Residual symmetry for $p=2$]
For $p = 2$, the closed form $\lambda_{2,x} = (\sqrt{x}-1)/(\sqrt{x}+1)$ yields the symmetry $|\lambda_{2,x}| = |\lambda_{2,1/x}|$: computing $\sqrt{x}$ and $\sqrt{1/x}$ by Pandrosion's method have exactly the same speed of convergence. For $p \geq 3$, this symmetry does not hold in general.
\end{remark}

%=====================================================================
\section{Optimal starting point}
\label{sec:optimal_start}
%=====================================================================

The number of iterations to reach precision $\delta$ is governed by $n_\delta \sim \ln(\eps_0/\delta) / \ln(1/\lambda_{p,x})$ (Proposition~\ref{prop:cout}). To minimize $n_\delta$, we must therefore minimize the initial residual $\eps_0 = |v_0 - \alpha|$.

\begin{proposition}[A canonical starting point]
\label{prop:optimal_start}
Among starting points of the form $s_0 = h^{(k)}(c)$, with $c$ a boundary value of the natural domain (i.e.\ $c \in \{0^+, 1^-\}$) and $k \geq 0$ a small integer, the choice
\begin{equation}
  \label{eq:optimal_s0}
  s_0^{\mathrm{opt}} := h(1) = 1 - \frac{x-1}{xp}
\end{equation}
achieves an initial gap to $s^*$ of order $|x - 1 - x\ln x|/(xp)$, which is $O(1/p)$ as $p \to \infty$ and strictly better than the naive choices $s_0 = 1/2$ or $s_0 = 1$ in the regime $p \gg 1$.
\end{proposition}

\begin{proof}
Since $s^* = x^{-1/p} = 1 - (\ln x)/p + O(1/p^2)$ and $h(1) = 1 - (x-1)/(xp)$,
\[
  |h(1) - s^*| \;=\; \left|\frac{x-1}{xp} - \frac{\ln x}{p}\right| + O(1/p^2)
  \;=\; \frac{|x - 1 - x\ln x|}{xp} + O(1/p^2),
\]
whereas $|1/2 - s^*| = |1/2 - 1 + (\ln x)/p + O(1/p^2)| = 1/2 + O(1/p)$ and $|1 - s^*| = (\ln x)/p + O(1/p^2)$. Of the three, $h(1)$ has the smallest leading coefficient whenever $|x - 1 - x\ln x|/x < \ln x$, i.e.\ essentially whenever $x > 1$ is not near $1$; the numerical gains are shown in Table~\ref{tab:optimal-start-expo}. The proposition does \emph{not} claim $h(1)$ is globally optimal; it claims only that it is a canonical, $\alpha$-independent choice with this specific quantitative guarantee.
\end{proof}

\begin{table}[h]
\centering
\begin{tabular}{cccccc}
\toprule
$(p,x)$ & $s^*$ & $s_0^{\mathrm{opt}} = h(1)$ & $|s_0^{\mathrm{opt}} - s^*|$ & Iters ($s_0^{\mathrm{opt}}$) & Iters ($s_0 = \tfrac{1}{2}$) \\
\midrule
$(2,2)$ & $0.7071$ & $0.7500$ & $0.043$ & $12$ & $13$ \\
$(3,2)$ & $0.7937$ & $0.8333$ & $0.040$ & $14$ & $16$ \\
$(4,2)$ & $0.8409$ & $0.8750$ & $0.034$ & $15$ & $17$ \\
$(3,8)$ & $0.5000$ & $0.7083$ & $0.208$ & $42$ & $0^*$ \\
\bottomrule
\end{tabular}
\caption{Comparison of iteration counts to reach $|v_n - \alpha| < 10^{-10}$. The canonical starting point $s_0^{\mathrm{opt}} = h(1)$ consistently reduces iterations. (*For $(3,8)$, $s_0 = 1/2$ happens to coincide with $s^*$.)}
\label{tab:optimal-start-expo}
\end{table}

%=====================================================================
\section{Steffensen acceleration: from linear to quadratic convergence}
\label{sec:steffensen}
%=====================================================================

Pandrosion's method converges linearly with ratio $\lambda_{p,x}$. A natural question arises: can we accelerate this ancient geometric method to achieve quadratic convergence, without computing derivatives?

\subsection{Aitken \texorpdfstring{$\Delta^2$}{Delta-squared} acceleration}

Given a linearly convergent sequence $(s_n)_{n \geq 0}$, Aitken's $\Delta^2$ process~\cite{stoer} defines the accelerated sequence
\[
  \hat{s}_n = s_n - \frac{(s_{n+1} - s_n)^2}{s_{n+2} - 2s_{n+1} + s_n}.
\]
When applied to Pandrosion's iterates, the sequence $(\hat{s}_n)$ converges \emph{faster} than $(s_n)$. Steffensen's method~\cite{steffensen} iterates this acceleration: at each step, it computes three consecutive Pandrosion iterates and applies Aitken's formula, then restarts from the accelerated value.

\begin{definition}[Steffensen--Pandrosion iteration]
\label{def:steffensen}
Define the composite map $T : s \mapsto \hat{s}$ by
\begin{equation}
  \label{eq:steffensen}
  T(s) = s - \frac{(h(s) - s)^2}{h(h(s)) - 2h(s) + s},
\end{equation}
where $h(s) = 1 - (x-1)/(x \cdot S_p(s))$ is Pandrosion's iteration.
\end{definition}

\begin{theorem}[Quadratic convergence of Steffensen--Pandrosion]
\label{thm:steffensen}
The Steffensen--Pandrosion iteration~\eqref{eq:steffensen} converges \emph{quadratically} to $s^* = x^{-1/p}$:
\[
  |T(s) - s^*| \leq K_S \cdot |s - s^*|^2 + O(|s - s^*|^3),
\]
where the Steffensen constant $K_S$ satisfies
\begin{equation}
  \label{eq:K_steffensen}
  K_S = \left|\frac{h''(s^*)}{2(1 - h'(s^*))}\right| = \left|\frac{h''(s^*)}{2(1 - \lambda_{p,x})}\right|.
\end{equation}
\end{theorem}

\begin{proof}
This is a standard result in fixed-point acceleration theory~\cite{stoer}. If $h$ has a fixed point $s^*$ with $h'(s^*) = \lambda \neq 1$ and $h$ is twice continuously differentiable near $s^*$, then the Steffensen map $T$ defined by~\eqref{eq:steffensen} satisfies $T(s^*) = s^*$ and $T'(s^*) = 0$, yielding quadratic convergence with the constant~\eqref{eq:K_steffensen}.
\end{proof}

\subsection{Numerical demonstration}

\begin{table}[H]
\centering
\begin{tabular}{cccc}
\toprule
Step $n$ & $s_n$ (Steffensen) & $|v_n - \cbrt{2}|$ & $|s_n - s^*|$ \\
\midrule
0 & $0.875000000000000$ & $2.71 \times 10^{-1}$ & $8.13 \times 10^{-2}$ \\
1 & $0.793949257650704$ & $7.90 \times 10^{-4}$ & $2.49 \times 10^{-4}$ \\
2 & $0.793700528604164$ & $8.32 \times 10^{-9}$ & $2.62 \times 10^{-9}$ \\
3 & $0.793700525984100$ & $< 10^{-15}$ & $< 10^{-16}$ \\
\bottomrule
\end{tabular}
\caption{Steffensen-accelerated Pandrosion for $p=3$, $x=2$. Machine precision is reached in \textbf{3~steps} (each step requires 2~evaluations of $h$, i.e.\ 6 evaluations of $S_p$ total). The step-$3$ output-residual is given as $< 10^{-15}$ rather than its raw floating-point value to avoid misreporting sub-ULP noise.}
\end{table}

\subsection{Comparison of quadratic constants}

Both Newton's method and Steffensen--Pandrosion converge quadratically, but with different constants.

\begin{table}[H]
\centering
\renewcommand{\arraystretch}{1.3}
\begin{tabular}{lcc}
\toprule
& \textbf{Newton} & \textbf{Steffensen--Pandrosion} \\
\midrule
Order & $2$ & $2$ \\
Constant ($p\!=\!3, x\!=\!2$) & $K_N = \dfrac{p-1}{2\alpha} \approx 0.794$ & $K_S \approx 0.013$ \\
Derivative required & Yes ($f'$) & No \\
Geometric origin & No & Yes \\
Steps to $10^{-15}$ & $5$ & $3$ \\
Evaluations of $S_p$ & --- & $6$ (total) \\
\bottomrule
\end{tabular}
\caption{Comparison of the two quadratically convergent methods for $\cbrt{2}$ on a per-iterate basis. Steffensen--Pandrosion has a much smaller quadratic constant ($K_S \ll K_N$), but requires two evaluations of $h$ per step against Newton's one derivative + one value evaluation. Per function evaluation the two methods are within a constant factor; the ``3 steps vs 5 steps'' figure is cosmetic and depends on the starting point.}
\end{table}

\begin{figure}[H]
  \centering
  \IfFileExists{pandrosion_steffensen.pdf}
    {\includegraphics[width=\textwidth]{pandrosion_steffensen.pdf}}
    {\fbox{\parbox{0.9\textwidth}{\centering \vspace{2em}\textit{Figure pandrosion\_steffensen.pdf not found.}\vspace{2em}}}}
  \caption{\textbf{Left:} Error decay comparison for $\cbrt{2}$, counted per iterate (not per function evaluation; see Remark~\ref{rem:fair-comparison}).
  Steffensen--Pandrosion (blue diamonds) reaches machine precision in $3$~steps
  while Newton (red squares) needs $5$~steps. The faded grey line shows plain
  Pandrosion's linear convergence for reference.
  \textbf{Right:} The quadratic constants $K_S \approx 0.013$ and
  $K_N \approx 0.794$ differ by a factor of ${\sim}\,61$. Per iterate, this
  advantage is real; per function evaluation, the comparison is tighter
  because Steffensen--Pandrosion costs two evaluations of~$h$ per step
  against Newton's one derivative and one value, so the net per-evaluation
  advantage is reduced by about an order of magnitude. The $\times 61$
  label refers to per-iterate convergence only.}
  \label{fig:steffensen}
\end{figure}

\begin{remark}
The construction illustrates, on a specific example, a standard point of acceleration theory: once an iteration converges linearly with a known ratio, Aitken $\Delta^2$ upgrades it to an order-$2$ scheme with a bounded constant. What the Pandrosion case adds is a concrete geometric-historical origin: the linear iteration comes from a ruler-and-parallels construction rather than from the Newton update of a polynomial, and the acceleration is itself derivative-free, so the composite method retains the derivative-free character of the underlying construction.
\end{remark}

%=====================================================================
\section{Scaling optimization: exploiting monotonicity}
\label{sec:scaling}
%=====================================================================

Theorem~\ref{thm:monotonicity} establishes that $\lambda_{p,x}$ is increasing in $x$: the larger the target $x$, the slower the convergence. For $x \gg 1$, the contraction ratio $\lambda_{p,x} \to 1^-$ (Proposition~\ref{prop:limits}), making direct application of Pandrosion's iteration impractical. We now show how to \emph{exploit} this very monotonicity to overcome it.

\subsection{The scaling principle}

The key observation is elementary: for any reference value $A > 1$ whose $p$-th root is known (or cheaply computable), we may write
\begin{equation}
  \label{eq:scaling}
  x^{1/p} = A^{1/p} \cdot \left(\frac{x}{A}\right)^{\!1/p}.
\end{equation}
Setting $x' = x/A$, Pandrosion's method is now applied to the \emph{reduced ratio} $x'$ instead of $x$. By choosing $A$ as the largest $p$-th power not exceeding $x$, we ensure $x' \in [1, (1+1/A^{1/p})^p)$, which is close to~$1$.

\begin{proposition}[Scaling reduction]
\label{prop:scaling}
Let $A = \lfloor x^{1/p} \rfloor^p$, so that $A^{1/p}$ is an integer and $x' = x/A \in [1, (1+1/\lfloor x^{1/p}\rfloor)^p)$. Then:
\begin{enumerate}
  \item The contraction ratio satisfies $\lambda_{p,x'} < \lambda_{p,x}$, and $\lambda_{p,x'} \to 0$ as $x \to \infty$ (so $x' \to 1^+$), while $\lambda_{p,x} \to 1^-$.
  \item The number of iterations drops from $O(\ln(\eps_0/\delta) / \ln(1/\lambda_{p,x}))$ to $O(\ln(\eps'_0/\delta) / \ln(1/\lambda_{p,x'}))$.
  \item The final result is recovered as $x^{1/p} = \lfloor x^{1/p} \rfloor \cdot (x')^{1/p}$.
\end{enumerate}
\end{proposition}

\begin{proof}
Part~(1) follows directly from Theorem~\ref{thm:monotonicity}: since $x' < x$ and $\lambda_{p,\cdot}$ is increasing, $\lambda_{p,x'} < \lambda_{p,x}$. Moreover, as $x \to \infty$, $x' \to 1^+$ so $\lambda_{p,x'} \to 0$ by Proposition~\ref{prop:limits}(1), while $\lambda_{p,x} \to 1^-$. Part~(2) follows from Proposition~\ref{prop:cout} applied to $x'$. Part~(3) is the factorization~\eqref{eq:scaling}.
\end{proof}

\begin{table}[h]
\centering
\begin{tabular}{ccccccc}
\toprule
$x$ & $\lambda_{3,x}$ & $A$ & $x' = x/A$ & $\lambda_{3,x'}$ & Reduction & Iterations \\
\midrule
$10$ & $0.615$ & $8$ & $1.25$ & $0.073$ & $8\times$ & $49 \to 8$ \\
$100$ & $0.890$ & $64$ & $1.5625$ & $0.145$ & $6\times$ & $213 \to 11$\\
$1000$ & $0.973$ & $729$ & $1.372$ & $0.103$ & $9\times$ & $>500 \to 9$ \\
$10000$ & $0.994$ & $9261$ & $1.080$ & $0.025$ & $39\times$ & $>500 \to 5$ \\
\bottomrule
\end{tabular}
\caption{Effect of scaling for $p = 3$ with $A = \lfloor x^{1/3} \rfloor^3$. The contraction ratio drops by up to $39\times$ and iteration counts collapse from hundreds to single digits. Iteration counts are for precision $10^{-10}$ with optimal starting point.}
\end{table}

\begin{remark}
The scaling reduction becomes more dramatic as $x$ grows: for $x = 10\,000$, the direct method requires over $500$ iterations while the scaled method needs only $5$. This is because $x' \to 1$ as $x \to \infty$, pushing $\lambda_{p,x'}$ toward the regime $\lambda \sim (p-1)(\alpha'-1)/2 \ll 1$ of Proposition~\ref{prop:limits}(1).
\end{remark}

\subsection{Geometric interpretation: Thales preconditioning}
\label{subsec:thales_preconditioning}

The scaling optimization admits an elegant geometric interpretation on Pandrosion's original Thales diagram.

\medskip

In the direct construction, one seeks to insert $p-1$ proportional means in a segment of ratio $x : 1$. When $x \gg 1$, this amounts to constructing parallels in a very elongated rectangle, where successive intercepts move slowly toward the target --- geometrically, the parallels are nearly horizontal.

\smallskip

The scaling strategy performs an \emph{homothety} on the construction: instead of starting from the unit segment, one starts from a segment of length $A^{1/p}$, chosen as the largest integer whose $p$-th power does not exceed $x$. The remaining construction now fills a much smaller gap: the rectangle has ratio $x' = x/A$ close to $1$, and the parallels converge rapidly.

\begin{proposition}[Double preconditioning]
\label{prop:preconditioning}
The scaling method combines two distinct acceleration mechanisms:
\begin{enumerate}
  \item \textbf{Initial residual reduction}: the reduced starting point $s_0^{\mathrm{opt}} = h_{x'}(1) = 1 - (x'-1)/(x'p)$ satisfies $|s_0^{\mathrm{opt}} - s'^*| \ll |s_0^{\mathrm{opt}} - s^*|$ since $x' \approx 1$ implies $s'^* = (x')^{-1/p} \approx 1$.
  \item \textbf{Intrinsic convergence acceleration}: each iteration contracts by $\lambda_{p,x'} \ll \lambda_{p,x}$, compounding the advantage over all subsequent steps.
\end{enumerate}
The total iteration count reduction is therefore superlinear in the ratio $\lambda_{p,x}/\lambda_{p,x'}$.
\end{proposition}

\begin{proof}
For mechanism~(1): $\eps'_0 = |v'_0 - \alpha'| \leq C \cdot |x' - 1| \to 0$ as $x' \to 1$, compared to $\eps_0 = |v_0 - \alpha| \geq C' \cdot |x - 1| \gg 1$ for the unscaled problem.

For mechanism~(2): the iteration count satisfies
\[
  n'_\delta \sim \frac{\ln(\eps'_0/\delta)}{\ln(1/\lambda_{p,x'})} \ll \frac{\ln(\eps_0/\delta)}{\ln(1/\lambda_{p,x})} \sim n_\delta,
\]
since both the numerator decreases (smaller $\eps'_0$) and the denominator increases (smaller $\lambda_{p,x'}$, hence larger $\ln(1/\lambda_{p,x'})$). These two effects multiply.
\end{proof}

\begin{remark}[From Pandrosion to a practical algorithm]
With scaling and Steffensen acceleration combined, Pandrosion's ancient geometric construction becomes a competitive root-finding algorithm: for any $x > 0$ and $p \geq 2$, one (i) reduces to $x' \approx 1$ via integer arithmetic, (ii) applies Steffensen--Pandrosion on $x'$ (converging in $2$--$3$ steps), and (iii) multiplies by $\lfloor x^{1/p} \rfloor$. The total cost is $4$--$6$ evaluations of $S_p$ plus one integer $p$-th root --- a purely algebraic, derivative-free computation whose origins lie in a ruler-and-parallels construction from fourth-century Alexandria.
\end{remark}

\subsection{Multi-precision benchmark against Newton and Halley}
\label{sec:benchmark}

We close Part~I with an experimental sanity check: how does the full Steffensen--Pandrosion--scaling pipeline compare against textbook Newton and Halley iterations for $x^{1/p}$ at high precision? The script \texttt{latex/benchmark\_multiprec.py} runs the three methods in MPFR-backed arithmetic (via the \texttt{mpmath} library) at precision targets of $50$, $200$, and $500$ decimal digits, starting all three methods from the same integer seed $\lfloor x^{1/p} \rfloor$ and stopping on the same condition $|u^p - x|/x \leq 10^{-\text{digits}}$. Table~\ref{tab:benchmark} summarises the iteration counts; wall-clock times are reported in the repository.

\begin{table}[H]
\centering
\small
\renewcommand{\arraystretch}{1.15}
\begin{tabular}{@{}lccccccccc@{}}
\toprule
& \multicolumn{3}{c}{\textbf{50 digits}} & \multicolumn{3}{c}{\textbf{200 digits}} & \multicolumn{3}{c}{\textbf{500 digits}} \\
\cmidrule(lr){2-4} \cmidrule(lr){5-7} \cmidrule(lr){8-10}
Regime & N & H & SP & N & H & SP & N & H & SP \\
\midrule
$p=2,\ x=2$             & 7 & 4 & 5 & 9  & 6 & 7 & 10 & 6 & 8 \\
$p=3,\ x=2$             & 7 & 4 & 5 & 9  & 6 & 7 & 10 & 6 & 8 \\
$p=3,\ x=11$            & 6 & 4 & 4 & 8  & 5 & 6 & 10 & 6 & 7 \\
$p=3,\ x=10^6+3$        & 4 & 2 & 2 & 6  & 4 & 4 & 7  & 5 & 5 \\
$p=5,\ x=99$            & 8 & 5 & 5 & 10 & 6 & 7 & 11 & 7 & 9 \\
$p=10,\ x=2^{20}+3$     & 4 & 2 & 2 & 6  & 4 & 4 & 7  & 5 & 5 \\
$p=20,\ x=2$            & 7 & 4 & 5 & 9  & 6 & 7 & 11 & 6 & 8 \\
\bottomrule
\end{tabular}
\caption{Iteration counts to reach the stated precision, starting from $u_0 = \lfloor x^{1/p} \rfloor$. \textbf{N} = Newton, \textbf{H} = Halley, \textbf{SP} = Steffensen--Pandrosion with scaling (Proposition~\ref{prop:scaling}). All three methods exhibit $\mathcal{O}(\log\log(1/\varepsilon))$ growth in iteration count; going from 50 to 500 decimal digits (a $10\times$ increase in precision exponent) adds about 3--4 iterations to each method, matching the theoretical $\log_2\log$ scaling.}
\label{tab:benchmark}
\end{table}

\paragraph{Observations.}
\begin{enumerate}
  \item \textbf{Iteration count scaling.} All three methods grow as $\log_2\log(1/\varepsilon)$, as predicted by theory (Newton and Steffensen--Pandrosion are order~$2$; Halley is order~$3$). In the table, going from 50 to 500 digits adds $\approx 3$--$4$ iterations, consistent with $\log_2 \log_{10}(10^{500}/10^{50}) = \log_2(500/50) \approx 3.3$.
  \item \textbf{Halley wins on count.} Halley (cubic order) consistently takes 2--3 fewer iterations than Newton (quadratic order), saving roughly $30\%$ on iteration count. Newton and Steffensen--Pandrosion--scaling differ by at most 1--2 iterations in every regime.
  \item \textbf{Wall clock.} Per-iteration, Newton is fastest (one power $u^{p-1}$ per step), Halley is close to Newton (one power $u^p$ and a rational combination), and Steffensen--Pandrosion--scaling is $2$--$3\times$ slower per iteration than Halley because each Aitken step requires two evaluations of the degree-$p$ polynomial $S_p$. On the iteration counts above and measured in the repository, Halley is the fastest method in total wall-clock time, Newton is within $10$--$30\%$, and Steffensen--Pandrosion--scaling sits $1.5$--$3\times$ slower.
  \item \textbf{The role of scaling.} Without scaling, Steffensen--Pandrosion on $x = 10^6 + 3$ with $p = 3$ would need $40$--$50$ iterations (the contraction ratio $\lambda_{3, 10^6+3} \approx 0.995$ gives $\log(10^{-200})/\log(1/0.995) \approx 92000$ iterations for the un-accelerated method, reduced by Aitken to $\sim 20$--$50$); with scaling, the count collapses to $4$. The scaling preconditioning is essential for Steffensen--Pandrosion to remain competitive at large $x$.
\end{enumerate}

\paragraph{Conclusion of the benchmark.} Steffensen--Pandrosion--scaling is not the fastest $p$-th root method in practice: Halley wins on both iteration count and wall clock, and Newton wins on wall clock by a small margin. What the benchmark does establish is that the scheme --- with scaling --- is \emph{competitive}: it achieves the same $\mathcal{O}(\log\log)$ asymptotic as Newton and Halley, within a small constant factor of them in every regime tested, while being derivative-free. For applications where the derivative is costly or unavailable, Steffensen--Pandrosion--scaling is a viable alternative; for standard $p$-th root computation, Halley remains the method of choice.

%=====================================================================
\part{Complex extension and formal verification}
%=====================================================================

Part~I developed the Pandrosion iterator $h_{p,x}(s) = 1 - (x-1)/(x S_p(s))$ on $\mathbb{R}_{>0}$, where $h_{p,x}$ has a unique fixed point $s^\ast = x^{-1/p}$ in $(0,1)$ and the Aitken--Steffensen upgrade $\sigma_{p,x}$ is a quadratic-contraction. On~$\mathbb{C}$ the situation is qualitatively different: $h_{p,x}$ has $p$ attractive fixed points (the $p$-th roots of $1/x$), separated by a non-trivial Julia set, and a classical impossibility result of McMullen~\cite{mcmullen} rules out purely iterative universal root finders for $p \geq 3$. Part~II extends the analysis to this complex setting via a multi-start scheme whose seed is target-dependent --- an $\alpha$-dependent preconditioning that takes the method out of the hypotheses of~\cite{mcmullen}.

\paragraph{Structural warning.} The scheme of Part II requires a \emph{target anchor} $\alpha$ (one of the $p$ solutions of $z^p = 1/x$) to be chosen before the seed is constructed. It is therefore a convergence scheme \emph{conditional on the target}, not a universal root-finder. Readers who expect the latter will misread Theorem~\ref{thm:multistart} below; the theorem says that, once a target is chosen, every starting point converges to that target, not that there is a parameter-free procedure picking out all roots. This is the unavoidable price of side-stepping McMullen's barrier.

\paragraph{From the real to the complex quadratic bound.} The Aitken--Steffensen map $\sigma_{p,x}$ is rational in $z$, so holomorphic on $\mathbb{C}$ away from the finite set $\mathrm{sing}(\sigma_{p,x}) := \{z : S_p(z) = 0 \text{ or } h_{p,x}(h_{p,x}(z)) - 2h_{p,x}(z) + z = 0\}$. At any cyclotomic fixed point $\alpha$ (i.e.\ $\alpha^p = 1/x$ with $S_p(\alpha) \neq 0$), the factorization of Part~I gives $\sigma_{p,x}(\alpha) = \alpha$ and $\sigma'_{p,x}(\alpha) = 0$.

We now cascade three quantities to avoid the apparent circularity between $K_\sigma$ and $R_\sigma$.

\medskip
\noindent\emph{Step 1 — fix a reference radius.} Set
\[
  R_0 \;:=\; \tfrac12\,\mathrm{dist}(\alpha, \mathrm{sing}(\sigma_{p,x})),
\]
so that $\overline{B}(\alpha, R_0) \subset \mathbb{C} \setminus \mathrm{sing}(\sigma_{p,x})$ and $\sigma_{p,x}$ is holomorphic on this closed disk.

\noindent\emph{Step 2 — explicit quadratic constant via Cauchy.} Define
\[
  M_0 \;:=\; \sup_{|w - \alpha| = R_0} |\sigma_{p,x}(w) - \alpha|,
  \qquad
  K_\sigma \;:=\; M_0 / R_0^2.
\]
By Cauchy's estimate applied to the holomorphic function $\sigma_{p,x}(\cdot) - \alpha$, which vanishes to order $2$ at $\alpha$ (since $\sigma_{p,x}(\alpha) = \alpha$ and $\sigma'_{p,x}(\alpha) = 0$), we obtain
\[
  |\sigma_{p,x}(z) - \alpha|
  \;\leq\; \frac{M_0}{R_0^2}\,|z - \alpha|^2
  \;=\; K_\sigma\,|z - \alpha|^2
  \qquad \text{for } z \in \overline{B}(\alpha, R_0).
\]

\noindent\emph{Step 3 — admissible contraction radius.} Define
\[
  R_\sigma \;=\; R_\sigma(x, p, \alpha) \;:=\; \min\!\bigl(\,R_0,\ 1/(2K_\sigma)\,\bigr).
\]
Then $R_\sigma \leq R_0$ (so the Cauchy bound of Step~2 is valid on $\overline{B}(\alpha, R_\sigma)$) and $K_\sigma R_\sigma \leq 1/2$. All three quantities $R_0$, $K_\sigma$, $R_\sigma$ are explicit functions of $(x, p, \alpha)$, computed in cascade without circularity. The seed construction below uses $K_\sigma$ and $R_\sigma$ as primitive quantities.

%=====================================================================
\section{Cyclotomic anchors and the complex iterator}
\label{sec:cyclotomic}
%=====================================================================

\begin{definition}[Cyclotomic anchors]\label{def:cyc-anchor}
For $\alpha \in \mathbb{C}^\ast$ and $p \geq 1$, let
\[
  \gamma_k := \alpha \cdot e^{2\pi i k / p}, \qquad k = 0, 1, \dots, p-1.
\]
\end{definition}

When $\alpha^p = 1/x$, the $\gamma_k$ are the $p$ distinct solutions of $z^p = 1/x$. By direct expansion of Definition~\ref{def:cyc-anchor}, $\omega \cdot \gamma_k = \gamma_{(k+1)\bmod p}$ (with $\omega = e^{2\pi i/p}$), $\sum_k \gamma_k = 0$, and $\prod_k \gamma_k = (-1)^{p-1} x$. These are definitional consequences for $z^p = x$, not specializations of the classical Galois or Vieta theorems; we use them freely.

The complex Pandrosion iterator is the analytic continuation of~$h_{p,x}$ to $\mathbb{C} \setminus \{z : S_p(z) = 0\}$, and the Aitken--Steffensen accelerator $\sigma_{p,x}$ of \S\ref{sec:steffensen} extends likewise. The quadratic local bound of Part~I applies verbatim to $\sigma_{p,x}$ in a disk $B(\gamma_k, R_\sigma)$ around every cyclotomic anchor, with the same explicit constants $K_\sigma$ and $R_\sigma$ satisfying $K_\sigma R_\sigma \leq 1/2$.

%=====================================================================
\section{Universal multi-start convergence in $\mathbb{C}$}
\label{sec:multistart}
%=====================================================================

\subsection{The multi-start seed}
\label{sec:seed}

\begin{definition}[Multi-start seed]\label{def:seed}
For a cyclotomic anchor $\alpha \in \mathbb{C}^\ast$ with $\alpha^p = 1/x$
and $z \in \mathbb{C}$,
\begin{equation}\label{eq:multistart-seed}
  \mathrm{seed}_{\alpha, \varepsilon}(z)
  \;:=\; \alpha + \varepsilon \cdot (z - \alpha),
  \qquad
  \varepsilon
  \;:=\; \frac{R_\sigma(x, p, \alpha)}{2\,\|z - \alpha\|}
  \quad (\varepsilon := 1 \text{ if } z = \alpha),
\end{equation}
where $R_\sigma(x, p, \alpha)$ is the explicit quadratic-contraction
radius from Part~I. The seed always lies in $\overline{B}(\alpha, R_\sigma/2)$.
\end{definition}

The seed plays in $\mathbb{C}$ the role that the canonical starting point of~\S\ref{sec:optimal_start} and the scaling of~\S\ref{sec:scaling} play on $\mathbb{R}$: it places the argument inside the quadratic-contraction disk of the chosen target before the iteration begins. It is $\alpha$-dependent by design, and this dependency is what moves the scheme out of the hypotheses of~\cite{mcmullen}.

\subsection{Universal multi-start theorem}

\begin{theorem}[Universal multi-start convergence]\label{thm:multistart}
Let $x \in \mathbb{C}^\ast$, $p \geq 1$, $\alpha \in \mathbb{C}^\ast$ with
$\alpha^p = 1/x$ and $S_p(\alpha) \neq 0$, and let $z \in \mathbb{C}$ with
$z \neq \alpha$. Define the seed parameter
\[
  \varepsilon_{\mathrm{seed}} \;:=\; \frac{R_\sigma(x,p,\alpha)}{2\,\|z - \alpha\|},
\]
which depends only on $(x, p, \alpha, z)$. Then for every target accuracy
$\delta > 0$, the integer
$N := N_{\mathrm{eff}}(K_\sigma, R_\sigma, \delta)$ of
Theorem~\ref{thm:eff-kt} satisfies
\[
  \forall n \geq N: \quad
  \big\|\sigma_{p,x}^n(\mathrm{seed}_{\alpha, \varepsilon_{\mathrm{seed}}}(z))
        - \alpha\big\| \;\leq\; \delta.
\]
The case $z = \alpha$ is trivial (take $\varepsilon_{\mathrm{seed}} := 1$ and
$N := 0$). The construction separates cleanly into two levels: the seed
parameter $\varepsilon_{\mathrm{seed}}$ and the admissible radius $R_\sigma$
are functions of $(x, p, \alpha, z)$ alone, while the iteration count $N$
alone depends on the target accuracy $\delta$.
\end{theorem}

\begin{proof}
Let $d := \|z - \alpha\|$; we may assume $d > 0$ (else $z = \alpha$ is already the target). With $\varepsilon_{\mathrm{seed}} := R_\sigma/(2d)$, set
\[
  z_0 := \mathrm{seed}_{\alpha, \varepsilon_{\mathrm{seed}}}(z) = \alpha + \varepsilon_{\mathrm{seed}}(z - \alpha).
\]
Then $\|z_0 - \alpha\| = \varepsilon_{\mathrm{seed}}\cdot d = R_\sigma/2 < R_\sigma$, so $z_0 \in B(\alpha, R_\sigma)$.

The quadratic-contraction bound of Part~I (whose proof applies verbatim to $\sigma_{p,x}$ viewed as an analytic map on any small enough disk centred at a fixed point, since $\sigma_{p,x}(\alpha) = \alpha$ and $\sigma_{p,x}'(\alpha) = 0$) reads: for $w \in B(\alpha, R_\sigma)$,
\[
  \|\sigma_{p,x}(w) - \alpha\| \leq K_\sigma \|w - \alpha\|^2,
\]
with $K_\sigma R_\sigma \leq 1/2$. Hence if $w_n := \sigma_{p,x}^n(z_0)$ satisfies $\|w_n - \alpha\| \leq R_\sigma$, then
$\|w_{n+1} - \alpha\| \leq K_\sigma \|w_n - \alpha\|^2 \leq K_\sigma R_\sigma \cdot \|w_n - \alpha\| \leq \tfrac12 \|w_n - \alpha\| \leq R_\sigma/2$,
so the orbit stays in $\overline{B}(\alpha, R_\sigma/2)$ by induction from $w_0 = z_0$.

Setting $e_n := \|w_n - \alpha\|$, one has $e_0 \leq R_\sigma/2$ and $e_{n+1} \leq K_\sigma e_n^2$. Theorem~\ref{thm:eff-kt} (proved below) then supplies the explicit $N = N_{\mathrm{eff}}(K_\sigma, R_\sigma, \varepsilon)$ with the stated asymptotics.
\end{proof}

\begin{remark}[The hypothesis $S_p(\alpha) \neq 0$]
The condition $S_p(\alpha) \neq 0$ rules out the degenerate case where the denominator of $h_{p,x}$ vanishes at the anchor. Since $S_p(\alpha) = (\alpha^p - 1)/(\alpha - 1)$ for $\alpha \neq 1$, one has $S_p(\alpha) = 0$ iff $\alpha^p = 1$ with $\alpha \neq 1$, which combined with $\alpha^p = 1/x$ forces $x = 1$ and $\alpha$ a non-trivial $p$-th root of unity. The hypothesis therefore fails on a single point of the $x$-line, and holds for every $x \neq 1$.
\end{remark}

\subsection{Explicit iteration count}

\begin{definition}[Effective iteration count]
For $K, R, \varepsilon > 0$ with $K R < 1$ and $K \varepsilon < 1$,
\[
  N_{\mathrm{eff}}(K, R, \varepsilon) :=
  \log_2 \big\lceil \log(K \varepsilon) / \log(K R) \big\rceil_+ + 1,
\]
with $\lceil \cdot \rceil_+$ the natural-number ceiling clamped at $0$.
\end{definition}

\begin{theorem}[Explicit iteration count]\label{thm:eff-kt}
For any quadratically contracting sequence $e_{n+1} \leq K \cdot e_n^2$
with $e_0 \leq R$, $K \cdot R < 1$, and $\varepsilon < R$, every iterate
$k \geq N_{\mathrm{eff}}(K, R, \varepsilon)$ satisfies $e_k \leq \varepsilon$.
Asymptotically $N_{\mathrm{eff}} \sim \log_2 \log(1/\varepsilon)$.
\end{theorem}

\begin{proof}
Set $u_n := K e_n$, so $u_0 \leq KR < 1$ and $u_{n+1} \leq u_n^2$. By induction $u_n \leq u_0^{2^n} \leq (KR)^{2^n}$. We want $u_n \leq K\varepsilon$, i.e.\ $(KR)^{2^n} \leq K\varepsilon$. Since $KR < 1$ and $K\varepsilon < 1$, both $\log(KR)$ and $\log(K\varepsilon)$ are negative; dividing by $\log(KR) < 0$ reverses the inequality direction and gives $2^n \geq \log(K\varepsilon)/\log(KR)$, i.e.\
\[
  n \;\geq\; \log_2 \Bigl\lceil \frac{\log(K\varepsilon)}{\log(KR)} \Bigr\rceil_+,
\]
which is achieved at $n = N_{\mathrm{eff}}(K, R, \varepsilon)$ by the $+1$ safety margin. The asymptotics $N_{\mathrm{eff}} \sim \log_2\log(1/\varepsilon)$ as $\varepsilon \to 0$ follow since $\log(K\varepsilon)/\log(KR) \sim \log(1/\varepsilon)/|\log(KR)|$.
\end{proof}

\subsection{Arbitrary target selector}

\begin{theorem}[Target selection]\label{thm:galois-equiv}
For every function $S : \mathbb{C} \to \mathrm{Fin}\,p$, every $z \in \mathbb{C}$, and every $\varepsilon > 0$, the multi-start scheme with target $\gamma_{S(z)}$ converges to $\gamma_{S(z)}$ with an iteration count of the form of Theorem~\ref{thm:eff-kt}.
\end{theorem}

Three selectors are formalized: constant (\texttt{principal\_selector}), Voronoï-nearest (\texttt{nearest\_anchor\_selector}), angular sector.

\subsection{The multi-start bundle theorem}
\label{sec:grand-master}

The previous subsections have presented the ingredients of the multi-start pipeline separately: the seed (Definition~\ref{def:seed}), the universal convergence theorem (Theorem~\ref{thm:multistart}), the target selectors (Theorem~\ref{thm:galois-equiv}), and the real-line preconditionings of Part~I transported to the cyclotomic anchor family. The following theorem bundles all five into a single statement. The Lean formalization refers to it as the ``multi-start grand master theorem'' (module \texttt{Foundations.multi\_start\_grand\_master}), a name we retain only to match the Lean naming convention; the theorem is a glue lemma, not a headline result.

\begin{theorem}[Multi-start bundle, ``grand master'' in Lean]\label{thm:grand-master}
Let $x \in \mathbb{C}^\ast$, $A \in \mathbb{C}^\ast$, $p \geq 1$, and $\alpha, \beta \in \mathbb{C}^\ast$ with $\alpha^p = x$ and $\beta^p = A$. Let $\gamma : \{0, 1, \ldots, p-1\} \to \mathbb{C}^\ast$ be a cyclotomic anchor family satisfying $\gamma_k^p = A/x$ and $S_p(\gamma_k) \neq 0$ for every $k$. Then for every query point $z_0 \in \mathbb{C}$, the following five properties hold simultaneously:
\begin{enumerate}
  \item[(1)] \textbf{Multi-start target selection.} There exists an index $s^\ast \in \{0, 1, \ldots, p-1\}$ --- the Voronoï-nearest anchor to $z_0$ --- with $\|\gamma_{s^\ast} - z_0\| \leq \|\gamma_t - z_0\|$ for every $t$.
  \item[(2)] \textbf{Steffensen fixing.} Every $\gamma_s$ is a fixed point of the scaled Steffensen map: $\sigma_{x/A,\,p}(\gamma_s) = \gamma_s$ for every $s$.
  \item[(3)] \textbf{Scaling compatibility.} The principal ratio satisfies $(\alpha/\beta)^p = x/A$.
  \item[(4)] \textbf{Reconstruction.} $\alpha = \beta \cdot (\alpha/\beta)$, so the original $p$-th root is recovered from the scaled one by multiplication by the integer root $\beta$.
  \item[(5)] \textbf{Optimal starting point.} The canonical seed for the scaled problem is $s_0^{\mathrm{opt}}(x/A) = h_{p,\,x/A}(1)$, matching Proposition~\ref{prop:optimal_start} of Part~I.
\end{enumerate}
\end{theorem}

\begin{proof}
(1) Voronoï-nearest existence: the function $t \mapsto \|\gamma_t - z_0\|$ attains its minimum on the finite set $\{0, \ldots, p-1\}$ at some index $s^\ast$.

(2) Steffensen fixing: by the factorization of Part~I specialized to $\gamma_s^p = A/x = 1/(x/A)$, $h_{p, x/A}(\gamma_s) = \gamma_s$; applying Aitken's $\Delta^2$ to a constant sequence $h(\gamma_s), h(h(\gamma_s)), \ldots$ yields $\sigma_{x/A, p}(\gamma_s) = \gamma_s$ whenever the denominator of the Aitken step is non-zero, which holds since $S_p(\gamma_s) \neq 0$ (the removable singularity at a fixed point is handled by the standard Aitken convention).

(3) Scaling compatibility: $(\alpha/\beta)^p = \alpha^p/\beta^p = x/A$ by direct computation.

(4) Reconstruction: $\beta \cdot (\alpha/\beta) = \alpha$ since $\beta \neq 0$.

(5) Optimal starting point: $h_{p, x/A}(1) = 1 - ((x/A) - 1)/((x/A) \cdot S_p(1)) = 1 - ((x/A) - 1)/((x/A) \cdot p)$, matching~\eqref{eq:optimal_s0} applied to the scaled problem $x' = x/A$.
\end{proof}

\begin{remark}[What the Grand Master theorem actually says]
Theorem~\ref{thm:grand-master} is not itself a new mathematical result; each of the five parts is proven elsewhere in Parts~I and~II. Its value is structural: it records, in a single statement, the \emph{compatibility} of the five layers of the pipeline --- selection, fixing, scaling, reconstruction, seeding --- so that downstream statements (in the formalization and in the effective iteration-count theorem below) can cite one name instead of five. The theorem's role is that of a \emph{glue lemma}: Proposition~(3), for instance, is used by the scaling preconditioning of Part~I to certify that a $\mathbb{R}$-level identity transports correctly to the cyclotomic family in $\mathbb{C}$.
\end{remark}

\begin{remark}[Effective version]
The effective counterpart of Theorem~\ref{thm:grand-master} replaces the asymptotic iteration count $N$ of Theorem~\ref{thm:multistart} with the explicit $N_{\mathrm{eff}}(K_\sigma, R_\sigma, \delta)$ of Theorem~\ref{thm:eff-kt}, applied per-anchor in the family $\gamma$. Concretely, for every $z_0 \in \mathbb{C}$ and every target accuracy $\delta > 0$, the following holds: choose $s^\ast$ as in Theorem~\ref{thm:grand-master}(1), form the seed $\mathrm{seed}_{\gamma_{s^\ast},\,\varepsilon_{\mathrm{seed}}}(z_0)$ with $\varepsilon_{\mathrm{seed}} = R_\sigma(x/A, p, \gamma_{s^\ast})/(2\|z_0 - \gamma_{s^\ast}\|)$, and iterate $\sigma_{x/A, p}$. After $N = N_{\mathrm{eff}}(K_\sigma, R_\sigma, \delta)$ steps, the iterate is within $\delta$ of $\gamma_{s^\ast}$, and multiplying by $\beta$ yields an approximation of $\alpha$ to comparable accuracy. The total cost is $\log_2 \log(1/\delta) + O(1)$ iterations plus one integer $p$-th root --- the scaling cost $\beta$ --- plus one final multiplication, exactly the pipeline the benchmark of \S\ref{sec:benchmark} exercises.
\end{remark}

\subsection{Explicit lower bound on the basin of attraction (unseeded iteration)}
\label{sec:basin-bound}

Theorem~\ref{thm:multistart} concerns the \emph{seeded} iteration: every starting point $z$ is first mapped via $\mathrm{seed}_{\alpha,\varepsilon_{\mathrm{seed}}}$ into the quadratic-contraction disk before $\sigma_{p,x}$ is applied. We now record a complementary, \emph{unconditional} statement about the \emph{unseeded} iteration: even without any preprocessing, there is an explicit disk around each cyclotomic anchor on which plain iteration of $\sigma_{p,x}$ already converges. This gives an explicit lower bound on the volume of the (true, dynamical) basin of attraction of each anchor.

\begin{definition}[Basin of attraction]
For a cyclotomic anchor $\gamma_k$ of $\sigma_{p,x}$, set
\[
  \Omega_k := \bigl\{ z \in \mathbb{C} :
    \sigma_{p,x}^n(z) \text{ is defined for all } n \geq 0
    \text{ and } \sigma_{p,x}^n(z) \to \gamma_k \text{ as } n \to \infty \bigr\}.
\]
\end{definition}

\begin{proposition}[Explicit basin lower bound]\label{prop:basin-lower-bound}
Under the standard hypotheses on $(x, p, \gamma_k)$ (namely $\gamma_k^p = 1/x$ and $S_p(\gamma_k) \neq 0$), the disk $\overline{B}(\gamma_k, R_\sigma(x, p, \gamma_k))$ lies entirely in $\Omega_k$. Consequently:
\begin{enumerate}
  \item[(i)] \textbf{Per-anchor volume:} $\mathrm{vol}(\Omega_k) \;\geq\; \pi\,R_\sigma(x, p, \gamma_k)^2$.
  \item[(ii)] \textbf{Disjointness:} the basins $\Omega_0, \ldots, \Omega_{p-1}$ are pairwise disjoint.
  \item[(iii)] \textbf{Total volume:} $\displaystyle \mathrm{vol}\!\bigl(\bigcup_{k=0}^{p-1}\Omega_k\bigr) \;\geq\; \sum_{k=0}^{p-1} \pi\,R_\sigma(x, p, \gamma_k)^2$.
\end{enumerate}
\end{proposition}

\begin{proof}
Fix $k$ and abbreviate $R := R_\sigma(x, p, \gamma_k)$, $K := K_\sigma(x, p, \gamma_k)$. From the three-step construction in Part~II's opening, we have $KR \leq 1/2$ and
\[
  |\sigma_{p,x}(z) - \gamma_k| \;\leq\; K\, |z - \gamma_k|^2
  \qquad \text{for all } z \in \overline{B}(\gamma_k, R_0),
\]
with $R \leq R_0$, so the bound holds a fortiori on $\overline{B}(\gamma_k, R)$.

\emph{Forward invariance.} For $z \in \overline{B}(\gamma_k, R)$,
\[
  |\sigma_{p,x}(z) - \gamma_k|
  \;\leq\; K\,|z - \gamma_k|^2
  \;\leq\; (K R)\,|z - \gamma_k|
  \;\leq\; \tfrac{1}{2}\,|z - \gamma_k|
  \;\leq\; R/2,
\]
hence $\sigma_{p,x}(\overline{B}(\gamma_k, R)) \subset \overline{B}(\gamma_k, R/2) \subset \overline{B}(\gamma_k, R)$.

\emph{Convergence.} Applying the inequality above inductively, $|\sigma_{p,x}^n(z) - \gamma_k| \leq 2^{-n}\,|z - \gamma_k| \leq 2^{-n} R \to 0$, so $z \in \Omega_k$. This proves the disk containment.

Part~(i) follows since $\overline{B}(\gamma_k, R) \subset \Omega_k$ gives $\mathrm{vol}(\Omega_k) \geq \pi R^2$.

Part~(ii) is immediate: if $z \in \Omega_j \cap \Omega_k$ with $j \neq k$, then $\sigma_{p,x}^n(z)$ would converge simultaneously to $\gamma_j$ and $\gamma_k$, which is impossible since $\gamma_j \neq \gamma_k$.

Part~(iii) follows from (i) and (ii) by additivity of Lebesgue measure.
\end{proof}

\begin{remark}[Comparison with Theorem~\ref{thm:multistart}]
Theorem~\ref{thm:multistart} covers every $z \in \mathbb{C}$ but only after pre-processing $z$ into a chosen target's disk via the seed. Proposition~\ref{prop:basin-lower-bound} does not need any pre-processing: the disk $\overline{B}(\gamma_k, R_\sigma)$ is directly in the unseeded basin. The two results are complementary: the seed (Theorem~\ref{thm:multistart}) addresses points outside any anchor's disk, while the basin bound (Proposition~\ref{prop:basin-lower-bound}) quantifies the unconditional contribution from points already close to an anchor.
\end{remark}

\begin{remark}[Tightness and numerical observations]
The bound $\pi R_\sigma^2$ is typically very loose. Numerical experiments (script \texttt{latex/basin\_estimate.py}) on $(p, x) = (3, 2)$ show that $\Omega_0$ (principal anchor) covers approximately $99.95\%$ of the grid $[-1.5, 1.5]^2$, while the analytic lower bound $\pi R_\sigma^2$ with an empirical $R_\sigma \geq 2$ gives $\mathrm{vol}(\Omega_0) \geq 4\pi \approx 12.57$. For $(p, x) = (7, 10)$ the radius tightens to $R_\sigma \approx 0.37$ empirically, showing that the geometric obstruction of a non-trivial Julia boundary begins to bite at larger $p$ and $x$. Proposition~\ref{prop:basin-lower-bound} is an explicit \emph{lower} bound only; improving it to a matching upper bound or a Hausdorff-dimension estimate on $\mathbb{C} \setminus \bigcup_k \Omega_k$ requires Fatou--Julia machinery not in Mathlib v4.7.0 and is left for future work.
\end{remark}

\begin{remark}[Pullback improvement]\label{rem:pullback}
The bound can be sharpened without leaving the elementary toolbox. Since $\sigma_{p,x}$ is a rational map, each iterate $\sigma_{p,x}^{-m}(\overline{B}(\gamma_k, R/2))$ is an open subset of $\Omega_k$ (pullback of a forward-invariant open set containing $\gamma_k$). Hence
\[
  \mathrm{vol}(\Omega_k)
  \;\geq\; \mathrm{vol}\!\Bigl(\bigcup_{m \geq 0} \sigma_{p,x}^{-m}\bigl(\overline{B}(\gamma_k, R/2)\bigr)\Bigr),
\]
and each pullback-level $m$ contributes a disjoint (up to measure-zero overlap) additional region. We do not pursue this refinement here because each pullback involves solving a degree-$d$ polynomial equation, and the bookkeeping obscures the elementary character of the $m = 0$ bound.
\end{remark}

\subsection{Multi-precision benchmark in the complex plane}
\label{sec:benchmark-complex}

Complementing the real-line benchmark of \S\ref{sec:benchmark}, the script \texttt{latex/benchmark\_complex.py} compares four methods on the complex problem $z^p = X$ for $X \in \mathbb{C}^\ast$:
\begin{itemize}
  \item Newton's iteration $z_{n+1} = ((p-1)z_n + X/z_n^{p-1})/p$;
  \item Halley's iteration $z_{n+1} = z_n \cdot ((p-1)z_n^p + (p+1)X)/((p+1)z_n^p + (p-1)X)$;
  \item Multi-start Steffensen--Pandrosion (Theorem~\ref{thm:multistart}), seeded at the principal anchor $\alpha = X^{1/p}$;
  \item Grand Master multi-start (Theorem~\ref{thm:grand-master}), combining scaling $\beta = \lfloor |X|^{1/p} \rfloor$, Voronoï-nearest selector on the anchor family, seed, and reconstruction.
\end{itemize}
All methods start from the same user-supplied $z_0 \in \mathbb{C}$ and stop when $|z_n^p - X|/|X| \leq 10^{-d}$ at precision target $d$ decimal digits. The seed radius is set uniformly to a conservative $R_\sigma = 0.5$.

\begin{table}[H]
\centering
\footnotesize
\renewcommand{\arraystretch}{1.15}
\begin{tabular}{@{}lcccccccc@{}}
\toprule
& \multicolumn{4}{c}{\textbf{50 digits}} & \multicolumn{4}{c}{\textbf{200 digits}} \\
\cmidrule(lr){2-5} \cmidrule(lr){6-9}
Regime & N & H & MS & GM & N & H & MS & GM \\
\midrule
$p=3,\ X=2+3i,\ z_0 = 1$                   & 8  & 5 & 7 & 7  & 10 & 6 & 9  & 9  \\
$p=3,\ X=10,\ z_0 = 2+i$                   & 8  & 5 & 7 & 6  & 10 & 6 & 9  & 8  \\
$p=4,\ X=-1+i,\ z_0 = 1+i$                 & 8  & 5 & 8 & 8  & 10 & 6 & 10 & 10 \\
$p=5,\ X=100+7i,\ z_0 = 2+i$               & 12 & 6 & 8 & 7  & 14 & 7 & 10 & 9  \\
$p=3,\ X=10^6+3,\ z_0 = 50$ (far)          & 9  & 5 & 6 & \textbf{4} & 11 & 6 & 8 & \textbf{6} \\
$p=3,\ X=10^6+17i,\ z_0 = 50+10i$          & 9  & 5 & 6 & \textbf{4} & 11 & 6 & 8 & \textbf{6} \\
$p=7,\ X=10,\ z_0 = 1+i$                   & 7  & 4 & 8 & 13 & 9  & 6 & 10 & 15 \\
\bottomrule
\end{tabular}
\caption{Complex-case iteration counts. \textbf{N} = Newton, \textbf{H} = Halley, \textbf{MS} = multi-start Pandrosion seeded at principal anchor, \textbf{GM} = Grand Master (scaling + Voronoï-nearest + seed). Stopping criterion $|z^p - X|/|X| \leq 10^{-d}$. Boldface marks the best (or tied-best) method in each row; Halley is best-or-tied in most regimes, Grand Master wins on large-$|X|$ regimes where scaling collapses the problem to $X' \approx 1$.}
\label{tab:benchmark-complex}
\end{table}

\paragraph{Observations.}
\begin{enumerate}
  \item \textbf{Halley is the iteration-count winner.} Cubic order gives Halley a consistent 30--40\% edge on iteration count in every tested regime. Wall-clock ordering is the same as the real case (Halley $<$ Newton $<$ MS $<$ GM).
  \item \textbf{Multi-start Pandrosion matches Newton within $\pm 2$ iterations.} The seed construction places $z_0$ inside the quadratic-contraction disk of the principal anchor, from which the $\log_2 \log$ scaling kicks in. The iteration count is within one or two of Newton's across all regimes.
  \item \textbf{Grand Master wins when scaling collapses the problem.} For $|X| \gtrsim 10^5$, Grand Master drops to $4$--$6$ iterations (at 50 and 200 digits respectively), a $30$--$40\%$ improvement over plain multi-start Pandrosion, because the scaled argument $X/\beta^p$ is close to $1$ and $\lambda_\infty(X/\beta^p)$ is small. This confirms experimentally the scaling preconditioning of Part~I, Proposition~\ref{prop:scaling}, in the complex setting.
  \item \textbf{Grand Master can be worse in niche regimes.} For $p = 7, X = 10, z_0 = 1+i$, Grand Master takes 13--15 iterations versus 8--10 for plain multi-start. The cause is that $|X|^{1/7} \approx 1.39 < 2$, so the integer scaling coefficient $\beta = 1$ provides no preconditioning, yet the Voronoï-nearest selector picks a non-principal anchor $\gamma_{s^\ast}$ for which the conservative $R_\sigma = 0.5$ is not tight, forcing extra pre-asymptotic iterations. This is a cost of the generality of the Grand Master pipeline.
  \item \textbf{All four methods achieve $\mathcal{O}(\log\log)$ scaling.} Going from $50$ to $200$ digits adds roughly $2$--$3$ iterations to each method, matching $\log_2(200/50) = 2$.
\end{enumerate}

\paragraph{Bottom-line (complex case).} As on the real line, Steffensen--Pandrosion-style methods are \emph{competitive} with textbook Newton/Halley at high precision: they remain within a small constant factor of Halley in every tested regime while being derivative-free, and the Grand Master pipeline provides an additional $30$--$40\%$ savings in the large-$|X|$ regime where integer scaling bites. For derivative-cheap settings, Halley remains the fastest; for derivative-expensive or derivative-unavailable settings (black-box polynomial evaluation), the Pandrosion pipeline is a viable alternative.

%=====================================================================
\section{Measurability and a.e.\ continuity of the solution map}
\label{sec:measurability}
%=====================================================================

The ``nearest-anchor'' target selector of Section~\ref{sec:multistart} uses the classical axiom of choice and is not, a priori, Borel-measurable. We record a canonical Borel version and the elementary analytic properties it enjoys.

\begin{definition}[Borel selector and Voronoï boundary]\label{def:borel-selector}
Given anchors $\gamma_0, \dots, \gamma_{p-1} \in \mathbb{C}$ (all distinct), define, for each $z \in \mathbb{C}$,
\[
  S_{\mathrm{Borel}}(z) := \min\bigl\{k \in \{0, \dots, p-1\} : \|\gamma_k - z\| = \min_{j} \|\gamma_j - z\|\bigr\},
\]
breaking ties by smallest index. The \emph{Voronoï boundary} of the anchor set is
\[
  \mathrm{VorBdy} := \bigl\{z \in \mathbb{C} : \exists\, j \neq k,\ \|\gamma_j - z\| = \|\gamma_k - z\|\bigr\}
  = \bigcup_{j < k} B_{jk},
\]
where $B_{jk} := \{z : \|\gamma_j - z\| = \|\gamma_k - z\|\}$ is the perpendicular bisector of $\gamma_j, \gamma_k$. Let $\mathrm{VorInt} := \mathbb{C} \setminus \mathrm{VorBdy}$.
\end{definition}

\begin{proposition}[Analytic structure of the solution map]
\label{thm:borel-combined}
The Borel solution map $\Phi : z \mapsto \gamma_{S_{\mathrm{Borel}}(z)}$ satisfies:
\begin{enumerate}
  \item[(i)] \textbf{Borel-measurability.} $\Phi$ is Borel-measurable $\mathbb{C} \to \mathbb{C}$.
  \item[(ii)] \textbf{Null boundary.} $\mathrm{VorBdy}$ has two-dimensional Lebesgue measure zero.
  \item[(iii)] \textbf{A.e.\ continuity.} $\Phi$ is locally constant on $\mathrm{VorInt}$, hence continuous there; combined with (ii), $\Phi$ is continuous at Lebesgue-a.e.\ $z \in \mathbb{C}$.
\end{enumerate}
\end{proposition}

\begin{proof}
\emph{(i) Borel-measurability.} The preimage $\{z : S_{\mathrm{Borel}}(z) = k\}$ is
\[
  \Bigl(\bigcap_{j \neq k} \{z : \|\gamma_k - z\| \leq \|\gamma_j - z\|\}\Bigr)
  \cap
  \Bigl(\bigcap_{j < k} \{z : \|\gamma_k - z\| < \|\gamma_j - z\|\}\Bigr),
\]
a finite Boolean combination of closed (resp.\ open) half-planes in $\mathbb{C}$, hence Borel. So $S_{\mathrm{Borel}}$ is Borel-measurable, and $\Phi$ is its composition with $k \mapsto \gamma_k$ (continuous, a fortiori Borel).

\emph{(ii) Null boundary.} Squaring $\|\gamma_j - z\| = \|\gamma_k - z\|$ and expanding in real coordinates yields the affine equation
\[
  B_{jk} = \bigl\{z \in \mathbb{C} : \mathrm{Re}\bigl(\overline{(\gamma_k - \gamma_j)}\,z\bigr) = \tfrac12(|\gamma_k|^2 - |\gamma_j|^2)\bigr\},
\]
a strict affine subspace of $\mathbb{C} \cong \mathbb{R}^2$ (Hausdorff dimension $1$). Each $B_{jk}$ therefore has two-dimensional Lebesgue measure zero. A finite union of null sets is null, so $\mathrm{vol}(\mathrm{VorBdy}) = 0$.

\emph{(iii) A.e.\ continuity.} Fix $z_0 \in \mathrm{VorInt}$ and let $k_0 := S_{\mathrm{Borel}}(z_0)$. By definition of $\mathrm{VorInt}$, $\|\gamma_{k_0} - z_0\| < \|\gamma_j - z_0\|$ for all $j \neq k_0$, so some $\delta > 0$ satisfies
$\|\gamma_{k_0} - z_0\| + 2\delta \leq \|\gamma_j - z_0\|$ for $j \neq k_0$. By continuity of $z \mapsto \|\gamma_j - z\|$, there is $\eta > 0$ such that on $B(z_0, \eta)$ both quantities change by at most $\delta$, giving $\|\gamma_{k_0} - z\| < \|\gamma_j - z\|$ for $j \neq k_0$, so $S_{\mathrm{Borel}}(z) = k_0$ and $\Phi(z) = \gamma_{k_0}$. Thus $\Phi$ is locally constant, hence continuous, on $\mathrm{VorInt}$. Combined with (ii), $\Phi$ is continuous at Lebesgue-a.e.\ $z$.
\end{proof}

For backwards reference we also expose the individual statements.

\begin{theorem}[Borel measurability]\label{thm:borel}
The map $S_{\mathrm{Borel}}$ is Borel-measurable; so is $\Phi$.
\end{theorem}

\begin{theorem}[Voronoï boundary is Lebesgue-null]\label{thm:julia-null}
$\mathrm{vol}(\mathrm{VorBdy}) = 0$.
\end{theorem}

\begin{theorem}[A.e.\ continuity]\label{thm:ae-continuous}
$\Phi$ is continuous at every $z \in \mathrm{VorInt}$; combined with Theorem~\ref{thm:julia-null}, $\Phi$ is continuous at Lebesgue-a.e.\ $z$.
\end{theorem}

Both theorems are direct consequences of Proposition~\ref{thm:borel-combined}, (ii) and (iii) respectively. Together with Theorem~\ref{thm:multistart}, they give the analytic character of the multi-start solution map: Borel-measurable everywhere, locally constant outside a Lebesgue-null set.

%=====================================================================
\section{Lean 4 formalization: corpus structure and elementary identities}
\label{sec:lean-corpus}
%=====================================================================

\subsection{Corpus overview}

The Lean development lives at
\url{github.com/ivan-fr/poussiere_de_x} and targets Lean 4 with Mathlib v4.7.0. It comprises 121 modules with zero \texttt{sorry}; every load-bearing theorem referenced in this paper has an audited axiomatic closure reducing to the Lean whitelist $\{\texttt{propext}, \texttt{Classical.choice}, \texttt{Quot.sound}\}$. Three layers of verifiability are offered to a third-party reviewer:
\begin{itemize}
  \item \textbf{Continuous integration.} A GitHub Actions workflow (\texttt{.github/workflows/lean-build.yml}) compiles the full corpus on every push; a green CI badge in the repository \texttt{README.md} certifies that the current \texttt{main} branch compiles cleanly, with no \texttt{sorry}.
  \item \textbf{Axiom audit file.} The file \texttt{AXIOMS.md} at the repository root contains the raw \texttt{\#print axioms} output for every load-bearing theorem cited in this paper; for each, the dependency closure reduces to the Lean whitelist $\{\texttt{propext}, \texttt{Classical.choice}, \texttt{Quot.sound}\}$.
  \item \textbf{Local reproduction.} Incremental build \texttt{bash scripts/lean-incremental.sh} takes 60--90~seconds on a modern laptop; per-theorem verification is available via \texttt{\#print axioms <theorem>} at the end of the relevant module.
\end{itemize}
The modules are organised in layers: foundations (geometric sums, fixed-point characterization, Voronoï basins, cyclotomic anchor definitions), core spine (explicit super-attractive rate, global $\log\log$ convergence, real-line $p = 2$ discharge), multi-start (universal convergence, target selection, iteration count, Grand Master theorem of \S\ref{sec:grand-master} via \texttt{Foundations.multi\_start\_grand\_master} and \texttt{PandrosionUltimateGrandMaster}), measure/topology (Borel measurability, Voronoï-null frontier, a.e.\ continuity), and a small set of classical specializations (Banach 1922, Aitken--Steffensen, Kronecker 1857, Aberth--Ehrlich, Kung--Traub attainability at $c=3$).

\paragraph{Statement-level \texttt{Prop}s that are NOT proven.} A single archive module \texttt{PandrosionOpenProblems.lean} contains named Lean \texttt{Prop}s whose proofs are \emph{deliberately absent}: the conjectural upper bound of Kung--Traub for $c \geq 3$, McMullen's 1987 impossibility theorem, Lehmer's Mahler-measure conjecture, and Smale's~17th problem. These are recorded as \texttt{axiom} or \texttt{Prop} declarations without accompanying proof terms. They are not used downstream by any load-bearing theorem of the corpus, and their presence in the repository should not be mistaken for formal verification of the corresponding results. A reader auditing the corpus can confirm this by checking that no \texttt{theorem} depending on these \texttt{Prop}s appears in the load-bearing chain. A future cleanup will move this archive module to a dedicated \texttt{archive} namespace to make the non-load-bearing status structurally explicit.

\subsection{Scope statement on classical open problems}

No claim of progress on classical open problems of polynomial root-finding is made. The corpus touches Kung--Traub ($c \geq 3$) through a compliance interface, not a proof, and records McMullen 1987 as a named statement without in-Lean proof (Mathlib v4.7.0 lacks the Fatou--Julia machinery needed). Lehmer's Mahler-measure conjecture and Smale's 17th problem are explicitly out of scope: the Pandrosion cyclotomic anchors produce algebraic integers whose Mahler measure is easy to compute (trivially far from Lehmer's extremal regime), and the multi-start $\mathcal{O}(\log\log(1/\varepsilon))$ complexity applies to the single-monomial family $z^p = x$, not to Smale's coefficient-uniform problem.

%=====================================================================
\part{Cyclotomic derivative-power sums and spectral balance}
%=====================================================================

Parts~I and~II developed the Pandrosion iterator and its multi-start extension. Part~III turns to a different object: the finite sum of inverse powers of the derivative $P'$ at the cyclotomic roots of $P(z) = z^p - 1/x$. These sums are finite, explicit, and admit vanishing identities whose structural content (``all integer arguments not divisible by $p$'') gives a uniform theorem from which the $\Sigma_P(1) = 0$, $Z_P(1) = 0$, and $Z_P(-1) = 0$ special values follow as corollaries. The part closes with a \emph{spectral balance} statement combining the vanishing identities with the a.e.-continuity and Julia-null results of Part~II.

All identities below have been verified symbolically in SymPy at precision $30$ decimal digits for $p \in \{2, 3, 5, 7, 10\}$ and representative $x$ (real and complex).

%=====================================================================
\section{The Pandrosion derivative-power sum $\Sigma_P$}
\label{sec:sigma-P}
%=====================================================================

Let $P(z) = z^p - 1/x$ with $x \in \mathbb{C}^\ast$, $p \geq 2$, and let $\alpha_0, \alpha_1, \ldots, \alpha_{p-1}$ denote the $p$ distinct roots of $P$, so that $\alpha_k^p = 1/x$ for every $k$. Fixing the principal root $\alpha := x^{-1/p}$ and $\omega := e^{2\pi i/p}$, the roots take the cyclotomic form $\alpha_k = \alpha \cdot \omega^k$. The derivative values are
\[
  P'(\alpha_k) \;=\; p\,\alpha_k^{p-1}
             \;=\; p \cdot \frac{\alpha_k^p}{\alpha_k}
             \;=\; \frac{p}{x\,\alpha_k},
\]
using $\alpha_k^p = 1/x$ in the last step. The multiplicative inverse is a simple expression:
\[
  \nu_k \;:=\; \frac{1}{P'(\alpha_k)} \;=\; \frac{x\,\alpha_k}{p}.
\]
We call $\nu_k$ the \emph{spectral weight} at the $k$-th root.

\begin{definition}[Pandrosion derivative-power sum]\label{def:sigma-P}
For $s \in \mathbb{Z}$, the Pandrosion derivative-power sum is
\[
  \Sigma_P(s) \;:=\; \sum_{k=0}^{p-1} P'(\alpha_k)^{-s}
             \;=\; \bigl(x/p\bigr)^{s} \sum_{k=0}^{p-1} \alpha_k^{s}.
\]
\end{definition}

The notation $\Sigma_P$ is chosen to avoid the letter~$\zeta$, which would spuriously suggest a Riemann-zeta connection. $\Sigma_P$ is a finite sum of $p$ explicit terms; the analogy with classical arithmetic zeta functions is typographic only.

%=====================================================================
\section{The general cyclotomic vanishing theorem}
\label{sec:cyclotomic-vanishing}
%=====================================================================

\begin{theorem}[General cyclotomic vanishing]\label{thm:cyclotomic-vanishing}
For every $p \geq 2$ and every $s \in \mathbb{Z}$ with $p \nmid s$,
\[
  \sum_{k=0}^{p-1} \alpha_k^s \;=\; 0.
\]
If instead $p \mid s$, then $\sum_{k=0}^{p-1} \alpha_k^s = p\,\alpha^s$.
\end{theorem}

\begin{proof}
Using $\alpha_k = \alpha \cdot \omega^k$,
\[
  \sum_{k=0}^{p-1} \alpha_k^{s}
   \;=\; \alpha^{s}\sum_{k=0}^{p-1} \omega^{ks}
   \;=\; \alpha^{s}\sum_{k=0}^{p-1} \bigl(\omega^{s}\bigr)^{k}.
\]
If $p \mid s$, then $\omega^{s} = 1$ and the sum equals $p$; otherwise $\omega^{s} \neq 1$ (it is also a $p$-th root of unity since $(\omega^{s})^{p} = (\omega^p)^s = 1$), so the geometric sum
\[
  \sum_{k=0}^{p-1}(\omega^s)^k \;=\; \frac{(\omega^s)^p - 1}{\omega^s - 1}
                                 \;=\; \frac{0}{\omega^s - 1} \;=\; 0.
\]
\end{proof}

\begin{corollary}[Vanishing of $\Sigma_P$]\label{cor:sigma-vanishing}
For every $p \geq 2$ and every $s \in \mathbb{Z}$ with $p \nmid s$,
$\Sigma_P(s) = 0$.
\end{corollary}

\begin{proof}
Immediate from Definition~\ref{def:sigma-P} and Theorem~\ref{thm:cyclotomic-vanishing}.
\end{proof}

\begin{corollary}[Closed-form product of derivative values]\label{cor:spectral-determinant}
For every $p \geq 1$ and $x \neq 0$,
\[
  \prod_{k=0}^{p-1} P'(\alpha_k)
   \;=\; (-1)^{p-1}\,\frac{p^p}{x^{p-1}}.
\]
\end{corollary}

\begin{proof}
From $P'(\alpha_k) = p/(x\alpha_k)$,
\[
  \prod_{k=0}^{p-1} P'(\alpha_k)
    \;=\; \frac{p^p}{x^p \cdot \prod_{k=0}^{p-1} \alpha_k}.
\]
Vieta's formulas applied to the monic polynomial $P(z) = z^p - 1/x$ (leading coefficient $1$, constant term $-1/x$) give $\prod_{k=0}^{p-1} \alpha_k = (-1)^p \cdot (-1/x) = (-1)^{p+1}/x$. Substituting,
\[
  \prod_{k=0}^{p-1} P'(\alpha_k)
    \;=\; \frac{p^p}{x^p} \cdot \frac{x}{(-1)^{p+1}}
    \;=\; (-1)^{p-1}\,\frac{p^p}{x^{p-1}},
\]
since $(-1)^{-(p+1)} = (-1)^{p+1} = (-1)^{p-1}$.
\end{proof}

%=====================================================================
\section{The energy-zeta $Z_P$ as $\Sigma_P(2\cdot)$}
\label{sec:energy-zeta}
%=====================================================================

Define the \emph{energy} at the $k$-th root by $E_P(\alpha_k) := P'(\alpha_k)^{2}$. The associated energy sum is
\[
  Z_P(s) \;:=\; \sum_{k=0}^{p-1} E_P(\alpha_k)^{-s}
        \;=\; \sum_{k=0}^{p-1} P'(\alpha_k)^{-2s}.
\]

\begin{theorem}[Energy-zeta reduction]\label{thm:energy-zeta-reduction}
For every $p \geq 1$, $x \neq 0$, and $s \in \mathbb{Z}$,
\[
  Z_P(s) \;=\; \Sigma_P(2s).
\]
\end{theorem}

\begin{proof}
$Z_P(s) = \sum_k (P'(\alpha_k)^2)^{-s} = \sum_k P'(\alpha_k)^{-2s} = \Sigma_P(2s)$, by definition.
\end{proof}

\begin{corollary}[Special values of $Z_P$]\label{cor:energy-special}
For $p \geq 3$:
\begin{enumerate}
  \item[(i)] $Z_P(0) = p$.
  \item[(ii)] $Z_P(1) = \Sigma_P(2) = 0$.
  \item[(iii)] $Z_P(-1) = \Sigma_P(-2) = 0$.
\end{enumerate}
\end{corollary}

\begin{proof}
For (i), $Z_P(0) = \sum_k 1 = p$. For (ii), apply Corollary~\ref{cor:sigma-vanishing} with $s = 2$: since $p \geq 3$ implies $p \nmid 2$, $\Sigma_P(2) = 0$. For (iii), the same corollary with $s = -2$: $p \geq 3$ implies $p \nmid -2$, so $\Sigma_P(-2) = 0$.
\end{proof}

\begin{remark}[Electrostatic interpretation]\label{rem:electrostatic}
Viewing the $p$ cyclotomic roots as unit charges placed at $\alpha_0, \ldots, \alpha_{p-1}$, with self-energy $E_P(\alpha_k) = P'(\alpha_k)^2$ and capacitance $1/E_P(\alpha_k)$, the three values above read:
\begin{itemize}
  \item total self-energy $\sum_k E_P(\alpha_k) = 0$ (for $p \geq 3$, total energy cancels cyclotomically);
  \item total capacitance $\sum_k 1/E_P(\alpha_k) = 0$ (same cancellation);
  \item only the count $Z_P(0) = p$ survives.
\end{itemize}
The Pandrosion cyclotomic anchor family is, in this sense, a \emph{perfectly balanced electrostatic system} for $p \geq 3$. For $p = 2$, the vanishings fail (since $p \mid 2$ in that case) and the system is degenerate.
\end{remark}

%=====================================================================
\section{Multi-start spectral balance}
\label{sec:spectral-balance}
%=====================================================================

Combining the vanishing identity of Corollary~\ref{cor:sigma-vanishing} with the a.e.-structural results of Part~II yields a unified statement.

\begin{theorem}[Multi-start spectral balance]\label{thm:spectral-balance}
Under the standard hypotheses of Part~II (namely $x \in \mathbb{C}^\ast$, $p \geq 2$, $\alpha \in \mathbb{C}^\ast$ with $\alpha^p = 1/x$, and $S_p(\alpha_k) \neq 0$ for every $k$), all five of the following hold simultaneously:
\begin{enumerate}
  \item[(A)] \textbf{Spectral-weight balance:} $\displaystyle \sum_{k=0}^{p-1} \nu_k = 0$, where $\nu_k = x\alpha_k/p$.
  \item[(B)] \textbf{Multi-start image:} for every $z \in \mathbb{C}$, the Borel multi-start target (Definition~\ref{def:borel-selector}) lies in $\{\alpha_0, \ldots, \alpha_{p-1}\}$.
  \item[(C)] \textbf{A.e.\ uniqueness:} for Lebesgue-a.e.\ $z \in \mathbb{C}$, the Voronoï-nearest anchor to $z$ is unique.
  \item[(D)] \textbf{Julia-null:} the Voronoï boundary $\mathrm{VorBdy}$ has 2-dimensional Lebesgue measure zero.
  \item[(E)] \textbf{A.e.\ continuity:} the Borel solution map $\Phi$ is continuous at Lebesgue-a.e.\ $z$.
\end{enumerate}
\end{theorem}

\begin{proof}
(A) is Corollary~\ref{cor:sigma-vanishing} at $s = 1$: $\Sigma_P(1) = 0$, i.e.\ $\sum_k 1/P'(\alpha_k) = \sum_k \nu_k = 0$. (B) is tautological from the definition of the Borel selector (Section~\ref{sec:measurability}). (C) follows from (D) since on the complement of the Voronoï boundary, the nearest-anchor map is single-valued. (D) is Theorem~\ref{thm:julia-null}. (E) is Theorem~\ref{thm:ae-continuous}.
\end{proof}

\begin{remark}[What the balance theorem says]
Theorem~\ref{thm:spectral-balance} is a glue statement: each of (A)--(E) is proven separately elsewhere. Its value is in the bundling: the cyclotomic-anchor family simultaneously satisfies a \emph{finite arithmetic constraint} (A), a \emph{measure-theoretic constraint} (D), and a \emph{regularity constraint} (E). This three-way compatibility justifies using cyclotomic anchors rather than an arbitrary root labelling.
\end{remark}

%=====================================================================
\part{Four exceptional theorems}
%=====================================================================

Part~IV collects four theorems that emerged from the formalization effort. Each is stated in ordinary mathematical prose with a self-contained proof using only the tools already developed (quadratic contraction, Vieta, cyclotomic symmetry, Taylor bound). No Fatou--Julia machinery or advanced complex dynamics is required.

%=====================================================================
\section{No-cycle theorem for multi-start Steffensen orbits}
\label{sec:no-cycle}
%=====================================================================

\begin{theorem}[No non-trivial cycles in the contraction disk]\label{thm:no-cycle}
Let $(x, p, \alpha)$ satisfy the standard hypotheses, and let $z_0 \in \overline{B}(\alpha, R_\sigma/2)$. If there exist integers $0 \leq n < m$ with $\sigma_{p,x}^n(z_0) = \sigma_{p,x}^m(z_0)$, then $\sigma_{p,x}^n(z_0) = \alpha$.
\end{theorem}

\begin{proof}
The three-step construction of Part~II supplies constants $K_\sigma, R_\sigma$ with $K_\sigma R_\sigma \leq 1/2$ and the quadratic bound
\[
  |\sigma_{p,x}(z) - \alpha| \;\leq\; K_\sigma\,|z - \alpha|^2
  \qquad \text{on } \overline{B}(\alpha, R_\sigma).
\]
For $z \in \overline{B}(\alpha, R_\sigma/2)$, i.e.\ $|z - \alpha| \leq R_\sigma/2$,
\[
  |\sigma_{p,x}(z) - \alpha|
   \;\leq\; K_\sigma\,|z - \alpha|\,\cdot\,|z - \alpha|
   \;\leq\; K_\sigma \cdot \frac{R_\sigma}{2} \cdot |z - \alpha|
   \;\leq\; \frac{1}{4}\,|z - \alpha|,
\]
using $K_\sigma R_\sigma \leq 1/2$. Hence $\sigma_{p,x}$ restricted to $\overline{B}(\alpha, R_\sigma/2)$ is $\tfrac14$-Lipschitz with $\alpha$ as fixed point, and $\sigma_{p,x}(\overline{B}(\alpha, R_\sigma/2)) \subset \overline{B}(\alpha, R_\sigma/8) \subset \overline{B}(\alpha, R_\sigma/2)$.

Write $w_j := \sigma_{p,x}^j(z_0)$. By induction on $j$, $|w_j - \alpha| \leq 4^{-j}|z_0 - \alpha|$. Suppose now $w_n = w_m$ with $0 \leq n < m$. By the same induction applied from $w_n$ onwards, $|w_m - \alpha| \leq 4^{-(m-n)}|w_n - \alpha|$. Combined with $w_n = w_m$, this gives
\[
  |w_n - \alpha| \;\leq\; 4^{-(m-n)}|w_n - \alpha|.
\]
Since $m - n \geq 1$, $4^{-(m-n)} \leq 1/4 < 1$, forcing $|w_n - \alpha| = 0$, i.e.\ $w_n = \alpha$. The orbit is constant at $\alpha$ from index $n$ onwards.
\end{proof}

\begin{corollary}[Acyclicity of seeded multi-start orbits]\label{cor:seeded-acyclic}
For every $(x, p, \alpha, z)$ in the hypotheses of Theorem~\ref{thm:multistart}, the orbit of $\mathrm{seed}_{\alpha, \varepsilon_{\mathrm{seed}}}(z)$ under $\sigma_{p,x}$ has no non-trivial cycles.
\end{corollary}

\begin{proof}
By construction the seed lies in $\overline{B}(\alpha, R_\sigma/2)$ (Definition~\ref{def:seed}), so Theorem~\ref{thm:no-cycle} applies.
\end{proof}

%=====================================================================
\section{The spectral determinant formula}
\label{sec:spectral-determinant-thm}
%=====================================================================

\begin{theorem}[Spectral determinant]\label{thm:spectral-det-main}
For every $p \geq 1$ and $x \in \mathbb{C}^\ast$,
\[
  \prod_{k=0}^{p-1} P'(\alpha_k)
   \;=\; (-1)^{p-1}\,\frac{p^p}{x^{p-1}}.
\]
\end{theorem}

\begin{proof}
This is Corollary~\ref{cor:spectral-determinant}, restated for visibility.
\end{proof}

\begin{corollary}[Discriminant of $z^p - 1/x$]\label{cor:discriminant}
The discriminant of the Pandrosion polynomial is
\[
  \mathrm{disc}(z^p - 1/x)
   \;=\; (-1)^{p(p+1)/2}\,\frac{p^p}{x^{p-1}}.
\]
\end{corollary}

\begin{proof}
For a monic polynomial $P$ of degree $p$, $\mathrm{disc}(P) = (-1)^{p(p-1)/2}\prod_k P'(\alpha_k)$. Combining with Theorem~\ref{thm:spectral-det-main}:
\[
  \mathrm{disc}(z^p - 1/x)
   \;=\; (-1)^{p(p-1)/2} \cdot (-1)^{p-1} \cdot \frac{p^p}{x^{p-1}}.
\]
The sign $(-1)^{p(p-1)/2 + (p-1)} = (-1)^{(p-1)(p/2 + 1)} = (-1)^{(p-1)(p+2)/2}$. A parity check (cases $p \equiv 0, 1, 2, 3 \pmod 4$) shows $(p-1)(p+2)/2 \equiv p(p+1)/2 \pmod 2$, so $\mathrm{disc} = (-1)^{p(p+1)/2}\,p^p/x^{p-1}$.
\end{proof}

Numerically verified for $(p, x) \in \{(3, 2), (5, 3), (7, 10), (3, 1{+}2i), (4, -1{+}i)\}$ to 30 decimal digits in arbitrary-precision arithmetic (\texttt{latex/benchmark\_complex.py}).

%=====================================================================
\section{Cyclotomic Voronoï equipartition}
\label{sec:voronoi-equipartition}
%=====================================================================

\begin{theorem}[Equipartition of cyclotomic Voronoï cells]\label{thm:voronoi-equipartition}
Let $\gamma_k = \alpha\omega^k$, $k = 0, \ldots, p-1$, be the cyclotomic anchor family, and let $V_k := \{z \in \mathbb{C} : \|\gamma_k - z\| < \|\gamma_j - z\|\ \forall j \neq k\}$ be the open Voronoï cell of $\gamma_k$. Then the rotation $T_\omega : z \mapsto \omega z$ satisfies $T_\omega(V_k) = V_{(k+1)\bmod p}$. Consequently, for every rotation-invariant Borel measure $\mu$ on $\mathbb{C}$ (e.g.\ Lebesgue measure restricted to any origin-centred disk),
\[
  \mu(V_0) \;=\; \mu(V_1) \;=\; \cdots \;=\; \mu(V_{p-1}).
\]
\end{theorem}

\begin{proof}
$T_\omega$ is an isometry of $\mathbb{C}$ (preserves Euclidean distance). It permutes anchors via $T_\omega(\gamma_k) = \omega \cdot \alpha\omega^k = \alpha\omega^{k+1} = \gamma_{(k+1) \bmod p}$. For $z \in V_k$, the Voronoï inequalities $\|\gamma_k - z\| < \|\gamma_j - z\|$ for $j \neq k$ become, after applying $T_\omega$, $\|\gamma_{k+1} - T_\omega z\| < \|\gamma_{j+1} - T_\omega z\|$ for $j + 1 \not\equiv k + 1 \pmod p$, which is exactly $T_\omega z \in V_{(k+1) \bmod p}$. Since $T_\omega$ is a bijection, $T_\omega(V_k) = V_{(k+1) \bmod p}$. Rotation-invariance of $\mu$ gives $\mu(V_{k+1}) = \mu(T_\omega V_k) = \mu(V_k)$, whence equipartition.
\end{proof}

\begin{corollary}[Per-cell lower bound]\label{cor:cell-bound}
Each Voronoï cell contains the disk of radius $\tfrac12 \min_{j \neq k}\|\gamma_j - \gamma_k\|$ centred at $\gamma_k$, hence $\mu(V_k) \geq \pi \cdot (\tfrac12 \min_{j\neq k}\|\gamma_j - \gamma_k\|)^2$ when $\mu$ is Lebesgue and the disk is contained in the support.
\end{corollary}

\begin{proof}
If $\|z - \gamma_k\| < \tfrac12 \min_{j\neq k}\|\gamma_j - \gamma_k\|$, then for any $j \neq k$, $\|z - \gamma_j\| \geq \|\gamma_j - \gamma_k\| - \|z - \gamma_k\| > \tfrac12 \|\gamma_j - \gamma_k\| \geq \|z - \gamma_k\|$, so $z \in V_k$.
\end{proof}

%=====================================================================
\section{Uniform log-log complexity over bounded $x$-regimes}
\label{sec:uniform-loglog}
%=====================================================================

\begin{theorem}[Uniform $\log\log$ complexity]\label{thm:uniform-loglog}
For every compact $K \subset (1, \infty)$ and every $\varepsilon \in (0, 1)$, there exists an integer $N(K, \varepsilon)$ such that for every $p \geq 2$, every $x \in K$, and every starting point $s_0 \in [1/2, 1]$, the Steffensen--Pandrosion iteration reaches precision $|v_n - x^{1/p}| \leq \varepsilon$ within at most $N(K, \varepsilon)$ steps. Furthermore, $N(K, \varepsilon) = \mathcal{O}(\log\log(1/\varepsilon))$ as $\varepsilon \to 0$, with constants depending only on $K$.
\end{theorem}

\begin{proof}
Set $\lambda^\ast(K) := \sup_{x \in K} \lambda_\infty(x) < 1$ (Proposition~\ref{prop:limits}(3): $\lambda_\infty$ is continuous on $(1, \infty)$ and strictly less than $1$ pointwise, so its supremum on the compact $K$ is strictly less than $1$). Fix $\delta := (1 - \lambda^\ast(K))/4 > 0$.

By the Taylor bound (inequality~\eqref{eq:taylor-bound}), for every $p \geq p_0(K) := \lceil \ln(\sup K) \rceil + 1$ and every $x \in K$,
\[
  \lambda_{p,x} \;\leq\; \lambda_\infty(x) + \frac{C(x)}{p}
              \;\leq\; \lambda^\ast(K) + \frac{C^\ast(K)}{p},
\]
where $C^\ast(K) := \sup_{x \in K} 3x(\ln x)^2/(x-1) < \infty$. Choose $p_1(K) := \max(p_0(K), \lceil C^\ast(K)/\delta \rceil)$, so that $\lambda_{p,x} \leq \lambda^\ast(K) + \delta \leq 1 - 3\delta$ for every $p \geq p_1(K)$ and $x \in K$.

For $2 \leq p < p_1(K)$, $\lambda_{p,x}$ is continuous on the compact $\{2, 3, \ldots, p_1(K)-1\} \times K$ and strictly less than $1$ pointwise (uniform monotonicity from Theorem~\ref{thm:monotonicity}); hence bounded away from $1$ by some $\lambda^{**}(K) < 1$. Set $\bar\lambda(K) := \max(1 - 3\delta, \lambda^{**}(K)) < 1$, so that $\lambda_{p,x} \leq \bar\lambda(K)$ for every $p \geq 2$, $x \in K$.

The un-accelerated Pandrosion iteration then satisfies $|s_n - s^\ast| \leq \bar\lambda(K)^n |s_0 - s^\ast|$; combined with the output-residual factor (Theorem~\ref{thm:nonasymp}) bounded on $K$, this gives linear convergence to $\alpha$ with uniform rate. The Aitken--Steffensen upgrade (Theorem~\ref{thm:steffensen}) turns this into quadratic convergence with constant $K_S(K) := \sup_{(p,x) \in \{p \geq 2\} \times K} |h''(s^\ast)/(2(1-\lambda_{p,x}))|$, uniform on $K$ by continuity plus $\lambda_{p,x} \leq \bar\lambda(K)$.

Combining: $n_1 := \lceil C_1(K) \log(1/R_\sigma(K)) \rceil$ linear iterations drive any $s_0 \in [1/2, 1]$ into the quadratic-contraction disk of radius $R_\sigma(K) := \inf_{(p,x)} R_\sigma(x,p,\alpha) > 0$; then $n_2 := N_{\mathrm{eff}}(K_S(K), R_\sigma(K), \varepsilon)$ Aitken--Steffensen iterations (Theorem~\ref{thm:eff-kt}) reach precision $\varepsilon$. The total $N(K, \varepsilon) = n_1 + n_2 = \mathcal{O}_K(\log\log(1/\varepsilon))$ is uniform in $p \geq 2$ and $x \in K$.
\end{proof}

This is the uniform version of the per-regime benchmark of~\S\ref{sec:benchmark}: the observed $\approx 3$--$4$ additional iterations when moving from 50 to 500 digits, \emph{across all tested $(p, x)$}, is a quantitative instance of Theorem~\ref{thm:uniform-loglog}.

%=====================================================================
\section{Conclusion}
\label{sec:conclusion}
%=====================================================================

The central tension of this paper is between Pandrosion's ancient geometric construction --- an order-one iteration requiring only a ruler and parallels --- and modern acceleration theory. Once the universal iteration $s_{n+1} = 1 - (x-1)/(x\,S_p(s_n))$ is written down, three ingredients suffice to make it competitive on~$\mathbb{R}$: the closed-form contraction ratio $\lambda_{p,x}$ (Theorem~\ref{thm:lambda}), the preconditionings of Propositions~\ref{prop:optimal_start} and~\ref{prop:scaling}, and the Aitken $\Delta^2$ upgrade (Theorem~\ref{thm:steffensen}). The ratio $\lambda_{p,x} = h'(s^*)$ is monotone in both $p$ and $x$ uniformly (Theorem~\ref{thm:monotonicity}, via the closed-form identities $P_p = p \cdot D_p'$ and $B_p = \sum_{j=1}^{p} j\alpha^{j-1}$), has an explicit large-$p$ limit $1 - \ln(x)/(x-1)$ (Proposition~\ref{prop:limits}), and admits a clean threshold analysis for $x < 1$ (Proposition~\ref{prop:x_small}).

On~$\mathbb{C}$, the extension via the target-dependent seed of Definition~\ref{def:seed} yields a universal multi-start convergence theorem (Theorem~\ref{thm:multistart}) with a closed-form $N_{\mathrm{eff}} \sim \log_2\log(1/\varepsilon)$ iteration count (Theorem~\ref{thm:eff-kt}), together with Borel-measurability and a.e.-continuity of the induced solution map (Theorems~\ref{thm:borel}--\ref{thm:ae-continuous}). All of these results are formally verified in the Lean~4 corpus described in Section~\ref{sec:lean-corpus}.

\medskip

The paper does not claim progress on classical open problems of polynomial root-finding, and does not claim a general algorithmic advantage over Newton or its modern variants. Its main deliverable is a concrete derivative-free root-finding scheme with closed-form convergence analysis and complete formal verification for the monomial family $z^p = x$, connecting a ruler-and-parallels construction from fourth-century Alexandria to contemporary formal-methods practice.

\paragraph{Perspectives.} We see four concrete, tractable extensions as immediate next steps:
\begin{enumerate}
  \item \textbf{Benchmark against Schr\"oder / K\"onig families of order $k \geq 3$} at multi-precision, to quantify how the factor-of-3 cubic gap of Halley over Newton compares against purely algebraic higher-order iterations (the König family is derivative-free in spirit and admits closed-form constants).
  \item \textbf{Extend to complex $x$} (not just complex anchor $\alpha$): the current Part~II assumes $x \in \mathbb{C}^\ast$ in the statements but the analysis is principally real for the scaling step. A fully complex scaling, via $\beta \in \mathbb{Z}[i]$ close to $X^{1/p}$ in the Gaussian integers, is a natural extension.
  \item \textbf{Formalize the non-asymptotic bounds of Section~\ref{sec:nonasymptotic} in Lean}, closing the gap between the analytic quantity $\Lambda$ and its in-Lean counterpart. This is a routine but useful extension of \texttt{UniformContractionRate.lean}.
  \item \textbf{Matching upper bound on the basin} (Remark~\ref{rem:pullback} pointed the way): show that the full basin $\Omega_k$ equals $\bigcup_{m \geq 0} \sigma_{p,x}^{-m}(\overline{B}(\gamma_k, R_\sigma))$ up to a Lebesgue-null set, yielding a matching upper bound on $\mathrm{vol}(\Omega_k)$. This stops short of requiring Fatou--Julia machinery.
\end{enumerate}
A Fatou--Julia formalization in Mathlib --- which would enable a proper in-Lean proof of McMullen's 1987 theorem --- is a much larger project and is mentioned here only for completeness, not as an immediate extension.

%=====================================================================
\begin{thebibliography}{9}

\bibitem{mactutor_pandrosion}
J.~J. O'Connor and E.~F. Robertson,
``Pandrosion of Alexandria,''
\textit{MacTutor History of Mathematics}, University of St~Andrews, 2020.
\url{https://mathshistory.st-andrews.ac.uk/Biographies/Pandrosion/}

\bibitem{mactutor_cube}
J.~J. O'Connor and E.~F. Robertson,
``Doubling the cube,''
\textit{MacTutor History of Mathematics}, University of St~Andrews, 2002.
\url{https://mathshistory.st-andrews.ac.uk/HistTopics/Doubling_the_cube/}

\bibitem{pappus}
Pappus of Alexandria,
\textit{Collectio}, Book~III (circa 340~AD).
Critical edition: F.~Hultsch, \textit{Pappi Alexandrini Collectionis quae supersunt},
Weidmann, Berlin, 1876--1878.

\bibitem{wantzel}
P.~Wantzel,
``Recherches sur les moyens de reconna\^itre si un probl\`eme de G\'eom\'etrie peut se r\'esoudre avec la r\`egle et le compas,''
\textit{Journal de Math\'ematiques Pures et Appliqu\'ees}, s\'er.~1, vol.~2, pp.~366--372, 1837.

\bibitem{stoer}
J.~Stoer and R.~Bulirsch,
\textit{Introduction to Numerical Analysis},
3rd~ed., Springer, New York, 2002.
Chap.~5: fixed-point iterations, convergence rates, and Aitken acceleration.

\bibitem{numrecipes}
W.~H. Press, S.~A. Teukolsky, W.~T. Vetterling, and B.~P. Flannery,
\textit{Numerical Recipes: The Art of Scientific Computing},
3rd~ed., Cambridge University Press, 2007.

\bibitem{steffensen}
J.~F. Steffensen,
``Remarks on iteration,''
\textit{Skandinavisk Aktuarietidskrift}, vol.~16, pp.~64--72, 1933.

\bibitem{aitken}
A.~C. Aitken,
``On Bernoulli's numerical solution of algebraic equations,''
\textit{Proceedings of the Royal Society of Edinburgh}, vol.~46, pp.~289--305, 1926.

\bibitem{ostrowski}
A.~M.~Ostrowski,
\textit{Solution of Equations and Systems of Equations},
2nd~ed., Academic Press, New York, 1966.

\bibitem{householder}
A.~S.~Householder,
\textit{The Numerical Treatment of a Single Nonlinear Equation},
McGraw-Hill, New York, 1970.

\bibitem{traub}
J.~F.~Traub,
\textit{Iterative Methods for the Solution of Equations},
Prentice-Hall, Englewood Cliffs NJ, 1964.

\bibitem{leanprover}
L.~de Moura and S.~Ullrich,
\textit{The Lean 4 Theorem Prover and Programming Language},
CADE-28, Springer LNCS, 2021.

\bibitem{mathlib}
The Mathlib Community,
\textit{Mathlib: A unified library of mathematics formalized},
\url{github.com/leanprover-community/mathlib4}.

\bibitem{kungtraub}
H.~T.~Kung and J.~F.~Traub,
``Optimal order of one-point and multipoint iteration,''
J.~ACM~\textbf{21} (1974), 643--651.

\bibitem{mcmullen}
C.~McMullen,
``Families of rational maps and iterative root-finding algorithms,''
Annals of Math.~\textbf{125} (1987), 467--493.

\end{thebibliography}

\end{document}
